% !TEX root = ../main.tex
\chapter{연구 방법}
\label{ch:methodology}

\ref{ch:literature}장은 다섯 가지 빈틈으로 끝났어요. 이 장은 그 빈틈마다 잴 수 있는 절차를 하나씩 정해 줘요. 또 이 논문의 방법인 결합 걷기 C-PR(결합 페이지랭크, 두 그래프를 함께 걷는 점수)과, 그 기준이 되는 한 층 방법 두 가지를 어떻게 평가하는지 밝혀요. 두 기준 방법은 엔도스랭크(EndorseRank)와 AWP(적응형 가중 페이지랭크, Adaptive Weighted PageRank)예요. 이 점수들은 아비트럼 원(Arbitrum One)의 스마트 계약(블록체인 위에서 저절로 실행되는 프로그램) 계정에 매기는 평판 점수예요. 엔도스랭크는 ERC-20(토큰 공통 규칙) 허락(allowance, 토큰을 얼마까지 대신 써도 된다는 허락) 그래프(점과 화살표로 이은 그림) 위에서 돌아가고, AWP는 송금(transfer) 그래프 위에서 돌아가요\parencite{do2023awp}. C-PR은 두 그래프를 하나의 걷기 안에서 함께 걸어요(\ref{sub:coupled-graph}절). 한 층만 걷는 두 걷기는 AWP를 만든 사람들이 정한 대로 화살표(엣지)에 무게(가중치)를 달아요. 그리고 모든 걷기는 무게를 단 페이지랭크(PageRank)로 풀어요\parencite{page1999pagerank}. AWP는 송금을 보낸 활동에 비례해서 걷기를 다시 출발시켜요(재시작). 엔도스랭크와 C-PR은 걷기를 고르게 다시 출발시켜요(균등 재시작).

방법 면에서 이 연구는 구성개념(점수가 재려고 하는 생각) 일치도, 곧 분야 일치도를 보는 방식을 따라요. 평판 점수는 이어지는 값으로 매긴 순위예요. 점검은 둘 중 하나로만 나뉘는 부도 라벨(빚을 갚지 못했는지 아닌지를 적은 정답 값) 대신, 계약 하나하나에서 볼 수 있는 통계에 바탕을 둬요. 순위를 매기는 방법 하나하나는, 직접 볼 수 없는 평판 구성개념을 가리키는 지표 하나예요. 이것은 구성 타당도(점수가 재려는 것을 정말 재는지)의 뜻을 따른 것이에요\parencite{cronbach1955construct}. 이 연구의 주장은 서로 견주는 주장이에요. 모든 점수를 각 구성개념의 점검과 견주고, 그 점수가 매끄럽게 다듬는(평활하는) 원시 차수(화살표를 그냥 센 수)와도 견주어요. 이 주장은 거래로 버는 돈이나 신용 부도(빚을 갚지 못하는 일)를 예측하는(미리 맞히는) 것이 아니에요.

실제 작업에서 내리는 선택은 모두 네 가지 원칙을 따라요.
\begin{itemize}
\item 같은 씨앗으로 견줄 수 있음: 모든 점수는 똑같은 계약 집합 위에서 계산해요. 그래서 점수 사이의 차이가 표본을 뽑는 방식 때문에 생길 수는 없어요.
\item 선택을 미리 정해 두기: 분석 계획과 버전을 매긴 설정 파일은 2026년 10월 9일 13:41~UTC에 공개 저장소에 커밋했어요(커밋 \texttt{43d8045}). 이 연구가 로그(계약이 남기는 기록)를 하나도 뽑기 전이었어요.\footnote{시각은 모두 UTC예요. 커밋에는 UTC+9로 기록되어 있어요.} 계획과 설정 파일에는 다음이 정해져 있어요. 기간, 코호트(묶음) 규칙과 계정 종류 규칙, 모든 점수와 그 매개변수, $\lambda$ 값, 대리 지표(대신 재는 값), W0 라벨, 부트스트랩(다시 뽑기), 벤치마크(속도 재기) 절차, 강건성 점검(조건을 바꿔도 결과가 버티는지 보기), 그리고 규칙~A예요. 규칙~A는 계획에 있는 단 하나의 판정 규칙이고, 그 대상은 C-PR이에요(연구 질문 RQ4; \ref{sec:holdout-protocol}절). 그 뒤로 스캔~A, 코호트, 에고 추출(코호트 계약에 닿는 기록만 뽑기), 같은 기간 결과와 W0 결과가 뒤따랐어요. 같은 날 16:49~UTC에는 W1의 등록 문서를 커밋했어요(커밋 \texttt{0d3a057}). 여기에는 허용 폭을 포함한 규칙~B, F4--F6, 민감도 분석(조건을 바꿔 결과가 얼마나 움직이는지 보기) 두 개, 그리고 두 화살표 종류 어느 것으로도 만들지 않은 결과에 대한 비교와 그 결론 문구가 들어 있어요(\ref{sub:fresh-holdout}절과 \ref{sub:neutral-label-protocol}절). 그다음에는 스캔~B, 등록한 평가, 보충 대비, 강건성 점검, 시빌 공격(가짜 주소를 많이 만들어 속이는 공격) 모형, 벤치마크가 이어졌어요. 이 결과들을 안 뒤에 수정 분석(등록 뒤에 금액을 달러로 바꿔 다시 한 분석)을 정의했어요(\ref{sec:amount-revision}절). 수정 분석은 금액을 달러(USD)로 매기고, 토큰 다섯 개만 쓰고, 점수를 각 동결일(점수를 고정하는 날) 전 12개월로 계산해요. \ref{ch:results}장이 보고하는 수정 분석은 같은 규칙을 이 점수들에 써요. 다른 대비는 모두 그 자료를 안 뒤에 계산했고, 차수 대비도 여기에 들어요. 이런 대비는 그렇다고 밝히고 구간과 함께 보고해요.
\item 만들기와 시간 재기: 각 코호트의 화살표 목록은 시간을 재기 시작하기 전에 한 번만 만들어요. 벤치마크는 풀이 프로그램(solver) 호출의 시간을 재요.
\item 범위: 가설 H1과, 주장 C4 가운데 순위 갈라짐을 다루는 절반은 엔도스랭크와 AWP의 줄 세우기를 견주어요. 나머지 가설들, 주장 C3과 C5, 그리고 C4 가운데 일치를 다루는 절반은 모든 점수에 관한 것이에요. C1과 C2는 결합 걷기와 그 빼 보기 실험에 관한 것이에요. C6은 트레이더(거래하는 사람) 결과에서의 허락한 주인 수(서로 다른 주인 수)와 송금 수, 그리고 그 페이지랭크들에 관한 것이에요. RQ3--RQ5는 W0과 W1의 스펜더(spender, 허락을 받아 대신 쓰는 쪽) 코호트를 쓰고, 규칙~B는 트레이더 코호트를 하나 더해요.
\end{itemize}

\dissertationSubheading[sec:notation]{기호}

이 장, 그리고 \ref{ch:results}장과 \ref{ch:discussion}장에서 쓰는 기호는 아래 표~\ref{tab:notation-methods}에 모아 두었어요. 줄인 목록은 표~\ref{tab:notation}에 있어요. 벡터(수를 줄지어 늘어놓은 것)는 굵은 글씨로, 집합은 필기체 글씨로 적어요.

\begingroup
\setstretch{1}\small\setlength{\tabcolsep}{3.5pt}\captionsetup{font=normalsize}
\begin{longtable}{lp{0.72\linewidth}}
\caption{그래프, 점수, 대리 지표, 매개변수의 기호.}
\label{tab:notation-methods}\\
\toprule
기호 & 뜻 \\
\midrule
\endfirsthead
\caption[]{(이어짐)}\\
\toprule
기호 & 뜻 \\
\midrule
\endhead
\bottomrule
\endlastfoot
$\mathcal{W}$ & 매칭 계약 코호트($|\mathcal{W}|=n=6{,}483$) \\
$G=(V,E)$ & 페이지랭크 풀이 프로그램에 넘기는 그래프: $V$는 남긴 화살표 $E$의 끝 점들 \\
$G_{\mathrm{approve}}=(V_A,E_A,w_A)$ & EndorseRank의 허락 그래프 \\
$G_{\mathrm{transfer}}=(V_T,E_T,w_T)$ & AWP의 송금 그래프 \\
$a^{\star}(u,v,\mathrm{token})$ & 한 토큰에서 주인 $u$가 스펜더 $v$에게 준, 기간 안의 마지막 허락(달러) \\
$\Delta t^{\star}(u,v,\mathrm{token})$ & $a^{\star}$ 값을 정한 허락의 나이 \\
$w_A(u,v)$ & EndorseRank의 화살표 무게, 주인 $u\to$ 스펜더 $v$: 토큰들에 걸쳐 더한 $\sigma(\Delta t^{\star})V(a^{\star})$ \\
$w_T(u,v)$ & AWP의 화살표 무게, 보낸 쪽 $u\to$ 받는 쪽 $v$: 송금들에 걸쳐 더한 $\sigma(\Delta t_e)V(x_e)$ \\
$x_e$ & 송금 $e$의 금액(달러) \\
$V(z)$, $b$ & AWP의 상한이 있는 값 변환, $V(z)=2/(1+e^{-bz})-1$, $b=\ln 3/m$(송금은 $b_T$, 허락은 $b_A$) \\
$m_T$, $m_A$ & W0 점수 기간에서 송금 금액의 중앙값($170$)과 허락 금액의 중앙값($100$), 달러 단위 \\
$U$ & C-PR과 받은 허락 금액에서 허락 하나에 두는 상한($19{,}000$달러) \\
$X_u$, $X^A_u$ & 송금 그래프와 허락 그래프에서 주소 $u$의 활발함(활동 정도), 곧 각각 AWP와 활동 재시작을 쓰는 EndorseRank의 다시 출발 무게 \\
$\tilde w_A(u,v)$, $\tilde w_T(u,v)$ & C-PR 두 층의 무게, 달러 단위 \\
$o(u)$; $o_A(u)$, $o_T(u)$ & 점 $u$의 나가는 무게의 합 $\sum_v w(u,v)$; 허락 층과 송금 층에서의 같은 값 \\
$P$ & 이동 행렬, $P_{u,v}=w(u,v)/o(u)$; 매달린 점의 줄은 모두 영 \\
$\mathbf{1}[\cdot]$ & 지시 함수: 조건이 맞으면 $1$, 아니면 $0$ \\
$\mathbf{1}_{\mathrm{dang}}$ & 매달린 점의 지시 벡터, 각 칸은 $\mathbf{1}[o(u)=0]$ \\
$\mathbf{r}$, $\mathbf{r}^{(t)}$ & 페이지랭크 벡터와 $t$걸음 뒤의 반복값 \\
$\boldsymbol{\pi}$ & 재시작(순간 이동) 분포, 매달린 점의 몫에도 씀: EndorseRank와 C-PR은 고르게, AWP는 $X_u$에 비례하게, 활동 재시작을 쓰는 EndorseRank는 $X^A_u$에 비례하게, S-PR은 허락 층 점수를 $V_T$로 좁힌 뒤 합이 하나가 되게 다시 맞춘 값 \\
$d$ & 페이지랭크 감쇠 계수(기본값 $0.85$) \\
$\varepsilon$ & 페이지랭크 수렴 허용 오차($10^{-8}$) \\
$I_{\max}$ & 페이지랭크 반복의 최대 횟수($300$) \\
$\lambda$ & 결합 연산자 C-PR에서 허락 층의 무게($\lambda=0.5$가 주된 값; $0.25$, $0.75$는 민감도; $\lambda=1$과 $\lambda=0$에서는 한 층만 걸음) \\
$\alpha_A(u)$, $\alpha_T(u)$ & 점 $u$의 C-PR 줄에서 두 층의 무게; $D_{\lambda}(u)$는 그 둘의 공통 분모 \\
$P^{(\lambda)}$ & C-PR의 결합 이동 행렬 \\
$\mathbf{s}_M$ & 방법 $M$의 평판 점수 벡터 \\
$\mathbf{p}_q$ & 대리 지표 $q$의 값 벡터 \\
$\sigma(\Delta t)$ & 허락이나 송금의 나이에 거는 로지스틱 감쇠 \\
$k$, $t_0$ & 로지스틱 감쇠의 가파름($0.01$)과 가운데 지점($180$일) \\
$\Delta t$ & 로지스틱 감쇠 기준 시점보다 며칠 전인지로 잰 이벤트의 나이($T_{\mathrm{obs}}$; W0에서는 $t_1$) \\
$T_{\mathrm{obs}}$ & 관측 기간이 끝나는 시각이자 W1의 동결일(2026-06-30 23:59:59 UTC) \\
$t_1$, $t_2$ & W0의 동결일(2026-03-31 23:59:59 UTC)과 그 라벨 기간의 끝(2026-06-30 23:59:59 UTC) \\
$\rho$ & 스피어먼의 순위 상관 \\
$R_i$, $S_i$ & 계약 $i$의 점수와 대리 지표 값의 중간 순위(같은 값들은 자리의 평균을 함께 가짐) \\
$\tau$, $\tau_b$ & 켄달의 순위 상관; 보고한 값은 모두 동점을 보정한 $\tau_b$ \\
$n_c$, $n_d$ & 순서가 맞는 계약 쌍과 어긋나는 계약 쌍의 수 \\
$n_0$; $n_1$, $n_2$ & 모든 쌍, $n(n-1)/2$; 첫 번째 순위와 두 번째 순위에서 동점인 쌍 \\
\end{longtable}
\endgroup

따로 말하지 않으면, 페이지랭크는 \textcite{page1999pagerank}의 페이지랭크에 무게를 단 형태를 뜻해요. 반복(한 바퀴) $t$에서 값을 새로 고치는 식은 아래와 같아요.
\begin{equation}
\label{eq:pagerank}
\mathbf{r}^{(t+1)}
=
(1-d)\,\boldsymbol{\pi}
+
d\,P^{\top}\mathbf{r}^{(t)}
+
d\,\bigl(\mathbf{r}^{(t)\top}\mathbf{1}_{\mathrm{dang}}\bigr)\,\boldsymbol{\pi},
\end{equation}
여기서 $P$는 아래 식~\eqref{eq:transition-impl}의 이동 행렬(전이 행렬)이에요. $\boldsymbol{\pi}$는 점들 위의 재시작 분포(어디서 다시 출발하는지의 비율)예요. 방법이 다른 분포를 정하지 않으면, 이 분포는 고르게($\pi_u=1/|V|$) 나뉘어요. 다른 분포를 정하는 방법은 \ref{sub:awp-graph}절의 AWP, \ref{sub:endorserank-graph}절의 활동 재시작(활발한 주소에서 더 자주 다시 출발하기)을 쓰는 엔도스랭크, \ref{sub:coupled-graph}절의 S-PR(보증 씨앗 페이지랭크)이에요. $\mathbf{1}_{\mathrm{dang}}$는 매달린 점(나가는 화살표가 없는 점)의 지시 벡터예요. 이 벡터의 각 칸은 $\mathbf{1}[o(u)=0]$이에요. 그래서 세 번째 항은 나가는 화살표가 없는 점들이 가진 몫을 모아요. $w(u,v)$를 $u$에서 $v$로 가는 방향이 있는 화살표의 무게라고 해요. 그리고 $o(u)=\sum_v w(u,v)$를 나가는 무게의 합이라고 해요. 그러면 이동 행렬은 아래와 같아요.
\begin{equation}
\label{eq:transition-impl}
P_{u,v}
=
\begin{cases}
w(u,v)/o(u) & \text{만약 }o(u)>0\text{이면},\\
0 & \text{그 밖의 경우.}
\end{cases}
\end{equation}
매달린 점의 줄은 모두 영이에요. 그래서 $P$에는 합이 하나가 되지 않는 줄이 있어요(이런 행렬을 아확률 행렬이라고 해요). 그 점들에서 사라진 몫은 식~\eqref{eq:pagerank}의 세 번째 항을 통해 돌아오고, $\boldsymbol{\pi}$에 따라 나뉘어요. 영인 줄마다 $1/|V|$를 채우면 식~\eqref{eq:transition}의 행 확률 행렬(줄마다 합이 하나인 행렬)이 돼요. 그리고 $\boldsymbol{\pi}$가 고르면, 식~\eqref{eq:pagerank}은 식~\eqref{eq:google-matrix}의 구글 행렬 형태와 같아져요. 반복은 $\mathbf{r}^{(0)}=\mathbf{1}/|V|$에서 시작해요. 그리고 $\|\mathbf{r}^{(t+1)}-\mathbf{r}^{(t)}\|_1 < \varepsilon$이 되거나 $I_{\max}$번 반복하면 멈춰요. $V$는 남긴 화살표의 끝 점들의 집합이에요. 마지막 벡터는 $V$ 위에서 합이 하나가 되도록 맞춰요. 계약은 점수가 큰 순서로도 늘어놓아요. 이때 동점은 주소로 순서를 정해요. 하지만 이 순번은 목록을 보여 줄 때만 써요. \ref{sec:statistical-tests}절의 상관(두 값이 함께 움직이는 정도)은 원래 점수 위에서 계산해요. 이 점수에서는 같은 점수가 그대로 동점으로 남아요.

감쇠 계수 $d$는 순수한 퍼뜨림(전파, $d\to 1$)과 고른 점수($d\to 0$) 사이를 이어요. 기본값 $0.85$는 웹에서 오래 써 온 설정이에요\parencite{page1999pagerank}. 강건성 점검에서 값을 바꿔 볼 때는 화살표 무게를 그대로 두고 이 값만 바꿔요(\ref{sec:scaling-robustness}절). 엄격한 $\varepsilon=10^{-8}$과 $I_{\max}=300$ 덕분에, 잰 실행 시간은 일찍 멈추는 일이 아니라 그래프의 구조에 따라 정해져요.

\dissertationSubheading[sec:data-sources]{자료 출처와 표본 뽑기}

이 절은 주 표본의 체인(블록체인), 기간, 공개 표, 이벤트, 계약을 정해요. 나중 단계 어디에서도 이 표본 밖의 이벤트를 들여오지 않아요. 그림~\ref{fig:data-overview}은 이 논문에서 쓴 온체인(블록체인 위의) 기록을 종류마다 하나의 달력 위에 놓아요. 그림에는 각 분석이 점수를 매기는 달, 라벨로 떼어 두는 달, 그리고 각 기록을 어디에 쓰는지도 함께 나와요.

\begin{figure}[htbp]
\centering
\begin{adjustbox}{max width=\linewidth}
% !TEX root = ../main.tex
% Blockchain data used in the thesis: one lane per analysis, one bar per kind of
% on-chain record, placed on the months it covers (1 Oct 2023 to 30 Sep 2026).
% The 18 months from October 2023 to March 2025 are drawn at a compressed scale
% (0.1 per month); the 18 months from April 2025 to September 2026 at 0.45 per month.
\begin{tikzpicture}[x=1cm, y=1cm, font=\scriptsize, oneline/.style={text height=1.6ex, text depth=0.45ex}]
\definecolor{dsApproval}{RGB}{0,114,178}
\definecolor{dsTransfer}{RGB}{0,158,115}
\definecolor{dsGMX}{RGB}{230,159,0}
\definecolor{dsOther}{RGB}{140,140,140}
\def\xa{1.95}% 1 Oct 2023
\def\xq{3.75}% 1 Apr 2025: start of the W0 score window and of the uncompressed scale
\def\xs{5.10}% 1 Jul 2025: start of the score window of T_obs
\def\xj{7.80}% 1 Jan 2026
\def\xw{9.15}% 1 Apr 2026: W0 freeze t1 (31 Mar 2026, 23:59:59 UTC)
\def\xt{10.50}% 1 Jul 2026: T_obs and W1 freeze (30 Jun 2026, 23:59:59 UTC)
\def\xe{11.85}% 1 Oct 2026: end of the W1 labels
\def\bh{0.36}% bar height
\def\xr{12.45}% left edge of the right-hand column
% \dsInput{colour}{y}{x start}{x end}{text}: records used as input
% (scores, cohorts, same-window checks).
\def\dsInput#1#2#3#4#5{%
  \fill[#1!35] (#3,{#2-\bh/2}) rectangle (#4,{#2+\bh/2});
  \draw[#1!75!black, line width=0.5pt] (#3,{#2-\bh/2}) rectangle (#4,{#2+\bh/2});
  \node[oneline, inner sep=0pt] at ({(#3+#4)/2},#2) {#5};}
% \dsHeld{colour}{y}{x start}{x end}{text}: labels counted after a freeze date
% and kept out of the scores.
\def\dsHeld#1#2#3#4#5{%
  \fill[pattern=north east lines, pattern color=#1] (#3,{#2-\bh/2}) rectangle (#4,{#2+\bh/2});
  \draw[#1!75!black, line width=0.7pt] (#3,{#2-\bh/2}) rectangle (#4,{#2+\bh/2});
  \node[oneline, fill=white, inner xsep=1pt, inner ysep=0pt] at ({(#3+#4)/2},#2) {#5};}
\def\rowlabel#1#2#3{\node[oneline, anchor=east, text=#1!75!black] at ({\xa-0.1},#2) {#3};}
\def\lanetitle#1#2{\node[oneline, anchor=west, font=\footnotesize\bfseries, fill=white, inner sep=1pt] at (0,#1) {#2};}
\def\laneuse#1#2{\node[anchor=north west, text width=4.9cm, align=flush left, inner sep=0pt] at (\xr,#1) {#2};}
\def\lanerule#1{\draw[black!25] (0,#1) -- (17.4,#1);}
\def\ytop{-0.3}% top of the lanes
\def\ybot{-10.5}% bottom of the lanes

% Calendar: year boundaries on the compressed part, month boundaries from April 2025
\foreach \x in {1.95,2.25,3.45} {\draw[black!12] (\x,\ytop) -- (\x,\ybot);}
\foreach \k in {0,...,18} {\draw[black!12] ({\xq+0.45*\k},\ytop) -- ({\xq+0.45*\k},\ybot);}
\foreach \k/\name in {0/4,1/5,2/6,3/7,4/8,5/9,6/10,7/11,8/12,9/1,10/2,11/3,12/4,13/5,14/6,15/7,16/8,17/9}
  {\node[oneline] at ({\xq+0.45*(\k+0.5)},0) {\name};}
\node[oneline] at ({(\xa+\xq)/2},0) {줄인 눈금};
\node[oneline, font=\scriptsize\bfseries] at ({(\xa+\xq)/2},0.34) {2023년 10월 -- 2025년 3월};
\node[oneline, font=\scriptsize\bfseries] at ({(\xq+\xj)/2},0.34) {2025년};
\node[oneline, font=\scriptsize\bfseries] at ({(\xj+\xe)/2},0.34) {2026년};
\draw[black!60] ({\xq-0.09},-0.24) -- ({\xq-0.01},-0.08);
\draw[black!60] ({\xq+0.01},-0.24) -- ({\xq+0.09},-0.08);
\draw[decorate, decoration={brace, amplitude=4pt}] ({\xa+0.03},0.62) -- ({\xt-0.03},0.62)
  node[oneline, midway, above=4pt] {스캔~A (등록 분석은 이 기간 전체로 점수를 매김)};
\draw[decorate, decoration={brace, amplitude=4pt}] ({\xt+0.03},0.62) -- ({\xe-0.03},0.62)
  node[oneline, midway, above=4pt] {스캔~B: W1 라벨};

% Freeze dates
\draw[black!70, dashed, line width=0.6pt] (\xw,\ytop) -- (\xw,\ybot);
\draw[black!70, dashed, line width=0.6pt] (\xt,\ytop) -- (\xt,\ybot);
\node[align=right, anchor=north east] at ({\xw+0.06},{\ybot-0.05}) {2026년 3월 31일\\W0 동결일 $t_1$};
\node[align=left, anchor=north west] at ({\xt-0.06},{\ybot-0.05}) {2026년 6월 30일\\$T_{\mathrm{obs}}$, W1 동결일};

% Lane 0: extraction
\lanetitle{-0.6}{추출}
\rowlabel{dsOther}{-1.0}{로그}
\dsInput{dsOther}{-1.0}{\xa}{\xt}{허락, 송금, GMX 로그 (스캔~A)}
\dsHeld{dsOther}{-1.0}{\xt}{\xe}{스캔~B}
\laneuse{-0.5}{계획의 코호트, 곧 어떤 토큰이든 허락을 받은 계약 $27{,}844$개의 에고 이벤트예요(\S\ref{sec:data-sources}).}
\lanerule{-1.35}

% Lane A: matched cohort selection
\lanetitle{-1.65}{매칭 코호트 고르기}
\rowlabel{dsApproval}{-2.05}{허락}
\dsInput{dsApproval}{-2.05}{\xs}{\xt}{다섯 토큰, 세 주인}
\laneuse{-1.55}{$T_{\mathrm{obs}}$ 전 열두 달 동안 적어도 세 주인에게서 다섯 토큰의 허락을 받은 계약이에요(\S\ref{sec:data-sources}).}
\lanerule{-2.40}

% Lane B: same-window analysis
\lanetitle{-2.70}{같은 기간 분석, 매칭 코호트($n=6{,}483$)}
\rowlabel{dsApproval}{-3.10}{허락}
\dsInput{dsApproval}{-3.10}{\xs}{\xt}{엔도스랭크; 허락 점검}
\rowlabel{dsTransfer}{-3.54}{송금}
\dsInput{dsTransfer}{-3.54}{\xs}{\xt}{AWP; 송금 점검과 시빌 점검}
\laneuse{-2.60}{두 종류의 화살표를 함께 쓰면 \mbox{C-PR}과 \mbox{S-PR}이 나와요. RQ1, 그리고 RQ2의 벤치마크와 RQ4의 같은 기간 부분이에요(\S\S\ref{sec:benchmark-results}--\ref{sec:rank-divergence}, \S\ref{sec:coupled-results}).}
\lanerule{-3.89}

% Lane C: holdout W0 on spender contracts
\lanetitle{-4.19}{홀드아웃 W0, 스펜더 계약($n=7{,}914$)}
\rowlabel{dsApproval}{-4.59}{허락}
\dsInput{dsApproval}{-4.59}{\xq}{\xw}{점수 입력; W0 코호트}
\dsHeld{dsApproval}{-4.59}{\xw}{\xt}{쌍}
\rowlabel{dsTransfer}{-5.03}{송금}
\dsInput{dsTransfer}{-5.03}{\xq}{\xw}{점수 입력}
\dsHeld{dsTransfer}{-5.03}{\xw}{\xt}{송금자}
\laneuse{-4.09}{$t_1$ 전 열두 달 안에 정해진 다섯 토큰의 허락을 가지고 있는 스펜더 계약이에요. RQ3과 RQ4예요(\S\ref{sec:holdout-protocol}, \S\S\ref{sec:holdout-results}--\ref{sec:cpr-results}).}
\lanerule{-5.38}

% Lane D: W0 trader run
\lanetitle{-5.68}{W0 트레이더 실행, 탐색적($n=87$)}
\rowlabel{dsApproval}{-6.08}{허락}
\dsInput{dsApproval}{-6.08}{\xq}{\xw}{점수 입력}
\rowlabel{dsTransfer}{-6.52}{송금}
\dsInput{dsTransfer}{-6.52}{\xq}{\xw}{점수 입력}
\rowlabel{dsGMX}{-6.96}{GMX V2}
\dsHeld{dsGMX}{-6.96}{\xw}{\xt}{결과}
\laneuse{-5.58}{$t_1$ 그래프에 있고 4--6월에 포지션을 세 번 이상 닫은 GMX~V2 계약 계정이에요. 이익 라벨에는 규칙이 없어요(\S\ref{sub:trader-holdout}). 청산 라벨은 결과 비교에만 들어가요.}
\lanerule{-7.31}

% Lane E: registered window W1
\lanetitle{-7.61}{등록 기간 W1}
\rowlabel{dsApproval}{-8.01}{허락}
\dsInput{dsApproval}{-8.01}{\xs}{\xt}{점수 입력}
\dsHeld{dsApproval}{-8.01}{\xt}{\xe}{쌍}
\rowlabel{dsTransfer}{-8.45}{송금}
\dsInput{dsTransfer}{-8.45}{\xs}{\xt}{점수 입력}
\dsHeld{dsTransfer}{-8.45}{\xt}{\xe}{송금자}
\rowlabel{dsGMX}{-8.89}{GMX V2}
\dsHeld{dsGMX}{-8.89}{\xt}{\xe}{결과}
\laneuse{-7.51}{스펜더 $7{,}960$개와 트레이더 $70$개이고, 모두 계약이에요. 고친 점수로 규칙~B와 결과 비교를 해요. 이 연구가 2026년 10월 9일 16:49~UTC의 등록 뒤에 뽑아냈어요(\S\ref{sub:fresh-holdout}).}
\lanerule{-9.24}

% Lane F: account type
\lanetitle{-9.54}{계정 종류}
\rowlabel{dsOther}{-10.0}{코드}
\fill[dsOther!75!black] (12.1,-10.0) +(0,0.13) -- +(0.13,0) -- +(0,-0.13) -- +(-0.13,0) -- cycle;
\node[anchor=east, fill=white, inner sep=1pt] at (11.9,-10.0) {고정한 블록에서 \texttt{eth\_getCode}, 2026년 10월 9일};
\laneuse{-9.44}{후보, 스펜더, 트레이더 주소 가운데 어느 것이 계약인지 알아봐요. EIP-7702 위임 표시는 EOA로 세어요(\S\ref{sec:data-sources}).}

% Legend: record colours, then bar styles
\begin{scope}[shift={(0,-11.6)}]
  \node[oneline, anchor=west, font=\scriptsize\bfseries, inner sep=0pt] at (0,0) {기록};
  \foreach \c/\t/\x/\y in {dsApproval/ERC-20 허락 로그/1.4/0, dsTransfer/ERC-20 송금 로그/5.0/0, dsGMX/GMX V2 포지션 닫기/8.6/0, dsOther/{모든 로그, 또는 RPC로 읽은 코드}/12.45/0} {
    \fill[\c!35] (\x,{\y-0.1}) rectangle ({\x+0.45},{\y+0.1});
    \draw[\c!75!black, line width=0.5pt] (\x,{\y-0.1}) rectangle ({\x+0.45},{\y+0.1});
    \node[oneline, anchor=west] at ({\x+0.5},\y) {\t};
  }
  \node[oneline, anchor=west, font=\scriptsize\bfseries, inner sep=0pt] at (0,-0.5) {막대};
  \fill[dsApproval!35] (1.4,-0.6) rectangle (1.625,-0.4);
  \fill[dsTransfer!35] (1.625,-0.6) rectangle (1.85,-0.4);
  \draw[black!60, line width=0.5pt] (1.4,-0.6) rectangle (1.85,-0.4);
  \node[oneline, anchor=west] at (1.9,-0.5) {칠한 막대: 점수, 코호트, 같은 기간 점검에 넣는 입력};
  \fill[pattern=north east lines, pattern color=black!70] (9.4,-0.6) rectangle (9.85,-0.4);
  \draw[black!60, line width=0.7pt] (9.4,-0.6) rectangle (9.85,-0.4);
  \node[oneline, anchor=west] at (9.9,-0.5) {빗금 막대: 떼어 둔 라벨, 동결일 뒤에 셈};
  \draw[black!70, dashed, line width=0.6pt] (1.625,-1.1) -- (1.625,-0.74);
  \node[oneline, anchor=west] at (1.9,-0.92) {점선: 동결일};
\end{scope}
\end{tikzpicture}

\end{adjustbox}
\caption[{이 논문에서 쓴 블록체인 자료: 출처, 다루는 달, 떼어 둔 라벨, 쓰임.}]{이 논문에서 쓴 블록체인 자료예요. 모두 아비트럼 원에서 나왔어요. 가로줄 하나는 분석 하나예요. 막대 하나는 온체인 기록 한 종류이고, 그 기록이 다루는 달 위에 놓여 있어요. 2025년 4월 전의 18개월은 줄여서 그렸어요. 색이 꽉 찬 막대는 점수, 코호트, 같은 기간 점검에 들어가는 입력이에요. 빗금 친 막대는 동결일 뒤에 센 라벨이고, 점수에는 넣지 않았어요. 이 연구는 7월부터 9월까지의 로그를 2026년 10월 9일 16:49~UTC에 W1의 등록 문서를 커밋한 뒤에야 뽑았어요. 막대 안의 글은 그 기록으로 무엇을 계산하는지 알려 줘요. 오른쪽 칸에는 코호트, 연구 질문, 결과를 보고하는 절이 적혀 있어요. 각 동결일 전 12개월 동안의 다섯 토큰 허락 로그로 그 동결일의 코호트를 골라요. 그다음 주 분석의 모든 허락 질의와 송금 질의를 이 코호트, 이 토큰들, 이 달들로 걸러요. 등록 분석(미리 등록한 그대로 한 분석)은 스캔~A의 모든 이벤트로 점수를 매겨요.}
\label{fig:data-overview}
\end{figure}

\begin{itemize}
\item 주 체인: 모든 이벤트는 아비트럼 원(체인 ID $42161$)에서 나와요. 아비트럼 원은 이더리움 위의 낙관적 롤업(일단 맞다고 보고, 문제가 있을 때만 따지는 롤업)이에요\parencite{thibault2022rollups}. 그 로그는 공개된 빅쿼리(BigQuery) 내보내기 자료에 들어 있어요. 체인을 하나만 쓰면 체인 사이의 별칭 문제를 피할 수 있어요. 별칭 문제란 한 사용자가 여러 체인에서 쓰는 주소들이 서로 다른 점으로 나타나는 일이에요. 또 표본의 모든 주소가 같은 수수료와 같은 지연 시간을 겪어요.

\item 추출 기간과 점수 기간: 스캔~A는 2023-10-01 00:00:00 UTC부터 2026-06-30 23:59:59 UTC까지의 이벤트를 뽑아요(질의에서는 2026-07-01을 끝 경계로 두고, 이 시각은 넣지 않아요). 스캔~B는 라벨 기간 W1, 곧 2026년 7월부터 9월까지의 이벤트를 뽑아요. 그래서 자료는 3년을 덮어요. 주 분석의 모든 점수는 그 동결일 전 12개월로 계산해요. 같은 기간 평가와 W1에서는 2025-07-01 00:00:00 UTC부터 $T_{\mathrm{obs}}=$ 2026-06-30 23:59:59 UTC까지이고, W0에서는 2025-04-01 00:00:00 UTC부터 $t_1=$ 2026-03-31 23:59:59 UTC까지예요(\ref{sec:holdout-protocol}절). 기간을 1년으로 두면, 점수, 대리 지표, 기준선(비교 기준), 라벨 이력에 들어가는 모든 이벤트의 나이에 상한이 생겨요. 또 두 동결일이 길이가 같은 기간을 가지게 돼요. 그리고 활동 기간 대리 지표와 활동한 달 수 대리 지표는 $365$일과 $12$개월까지 갈 수 있어요. 두 그래프는 똑같은 날짜를 써요. 로지스틱 감쇠(오래된 일일수록 가볍게 세는 무게)의 기준 시점은 동결일이에요. 그래서 마지막 날의 이벤트는 $0\le\Delta t<1$이고, 첫날의 이벤트는 나이가 $365$일이에요. 등록 분석(부록~\ref{sec:map-registered})은 2023-10-01부터 스캔~A의 모든 이벤트로 점수를 매겨요.

\item 빅쿼리 출처: 로그는 아래의 공개 빅쿼리 표에서 읽어요.
\begin{center}
\texttt{bigquery-public-data.goog\_blockchain\_arbitrum\_one\_us.logs}.
\end{center}
기록 하나는 계약이 내보낸 이벤트 하나예요. 기록에는 블록 타임스탬프(블록이 만들어진 시각), 주소, 토픽(이벤트의 머리 정보), 데이터가 들어 있어요. 스캔~A와 스캔~B는 이 연구에 필요한 이벤트를 이 연구의 구글 클라우드(Google Cloud) 프로젝트로 옮겨 적어요. 그 뒤의 질의는 이 옮겨 적은 것을 읽어요(\ref{sec:collection-strategy}절). 구글 클라우드 계정이 있는 사람은 누구나 같은 공개 표에 같은 질의를 다시 돌려 볼 수 있어요. 부록~\ref{ch:appendix-sql}에 질의의 모양이 나와 있어요.

\item 평판 이벤트: 엔도스랭크는 ERC-20 \texttt{Approval} 로그를 써요\parencite{vogelsteller2015erc20}. 이 로그에는 EIP-2612 \texttt{permit}(서명 허락)으로 한 허락도 들어 있어요\parencite{lundfall2020permit}. AWP는 \texttt{Transfer} 로그를 써요. 두 로그 모두 토픽이 정확히 세 개인 로그만, 그리고 \ref{sec:amount-revision}절의 다섯 토큰의 로그만 남겨요. 트레이더 라벨은 GMX~V2에서 포지션(거래하며 들고 있는 몫)을 닫은 기록에서 나와요.

\item 점수를 매기는 주소: 점수는 계약 계정에만 매겨요. 외부 소유 계정(EOA, 개인 키로 움직이는 계정)은 사용자 한 사람의 지갑이에요. EOA는 거래, 보관, 지불, 브리지(체인 사이를 잇는 다리)로 옮기기, 에어드롭 파밍(공짜 토큰을 노린 가짜 활동), 자동으로 움직이는 봇에 쓰일 수 있어요. 그리고 사용자마다 이 쓰임들을 섞는 비율이 서로 달라요. 그래서 EOA에는 한 줄로 세울 바탕이 되는 공통된 활동 모양이 없어요. 계약의 역할은 그 코드가 정해요(라우터, 금고(vault), 풀(pool) 같은 역할이에요). 그리고 한 계약에 허락을 해 주는 주인(허락을 해 주는 쪽)들은 그 역할에 맞는 같은 종류의 권한을 줘요. 그러니 서로 다른 계약이 받는 허락과 송금을 서로 견줄 수 있어요. 이런 까닭으로 EOA는 주인, 보낸 쪽(송신자), 받는 쪽(수신자)으로만 그래프에 들어가요. 하지만 이렇게 견줄 수 있는 것은 한 역할 안에서예요. 코호트에는 여러 역할이 섞여 있고, 개인이 지갑처럼 쓰는 계약도 들어 있을 수 있어요. 그리고 어떤 분석도 역할로 묶음을 나누지(층화하지) 않아요. 그래서 역할 사이의 차이가 모든 계수에 섞여 들어가요. 블록 $513{,}216{,}506$(2026년 10월 9일)에서 \texttt{eth\_getCode}가 EIP-7702 위임 표시가 아닌 코드를 돌려주면, 그 주소를 계약으로 쳐요. EIP-7702 위임 표시는 EOA가 자기 키를 그대로 가진 채 붙일 수 있는 표시예요\parencite{buterin2024eip7702}. 코호트 후보 밖의 주소는 세 시간 안의 더 나중 블록에서 읽었어요(부록~\ref{sec:account-type-classification}).

\item 매칭 코호트(모든 점수가 똑같이 점수를 매기는 계약 묶음): 일치도는 $n=6{,}483$개 계약에서 재요. 이 계약들은 $T_{\mathrm{obs}}$의 점수 기간 안에, 서로 다른 주인 적어도 세 명에게서 다섯 토큰 가운데 하나의 0이 아닌 \texttt{Approval}을 받은 계약이에요. 곧 계획의 코호트($27{,}844$개 계약, 2023-10-01부터 어떤 토큰으로든 허락을 받은 계약) 가운데, 그 기간 안에 다섯 토큰으로 이 규칙을 채우는 부분이에요. W0은 같은 규칙으로 자기 점수 기간 안에서 자기 코호트를 따로 뽑아요. 그래서 $t_1$ 뒤의 어떤 이벤트도, W0 자료를 모을 때 중심이 되는 계약을 고르는 데 끼지 않아요(\ref{sec:holdout-protocol}절). 이 목록은 허락과 코드 점검만으로 정해요. 하지만 허락 그래프와 상관없는 목록은 아니에요. 코호트 계약은 모두 기간 안 어느 때인가 허락을 가지고 있었기 때문이에요. 코호트 계약이 어느 그래프에든 꼭 나와야 하는 것은 아니에요. 어떤 그래프에 없으면, 그 그래프를 쓰는 방법은 그 계약에 $0$점을 줘요(\ref{sec:graph-construction}절). 이런 계약의 수는 \ref{sub:cohorts}절에서 세어요.

\item 에고 네트워크(코호트 계약에 닿는 화살표만 모은 그래프): 코호트 계약이 주인이나 스펜더인 이벤트, 또는 0이 아닌 송금의 보낸 쪽이나 받는 쪽인 이벤트만 써요. 추출에서는 계획의 코호트에 대한 이런 이벤트를 남겨요. 그리고 각 동결일에서는 자기 코호트에 대한 이런 이벤트를 자기 점수 기간과 라벨 기간에서 남겨요. 이렇게 하면 뽑는 자료의 양에 상한이 생겨요. 또 코호트 밖의 두 주소 사이의 이벤트는 한 번도 쓰지 않으므로, 각 그래프는 코호트를 둘러싼 한 단계 에고 네트워크가 돼요. 그래서 코호트 밖의 주소는 코호트 계약과 주고받은 이벤트를 통해서만 보여요. 허락한 주인이 세 명보다 적은 홀드아웃(나중 기간을 떼어 두고 맞혀 보는 시험) 스펜더나, 거의 모든 GMX 트레이더가 그런 주소예요. 이런 주소의 차수, 활발함, 점수, 그리고 스펜더의 라벨은 모두 그런 이벤트만 세어요. 보증(믿고 맡김)은 화살표 하나하나가 코호트 계약에 닿는 길을 따라서만 한 단계보다 멀리 퍼져요. 허락 걷기가 얼마나 그렇게 퍼지는지는 \ref{sub:theory-closed-form}절에서 살펴봐요.

\item 빼는 것: 멤풀(아직 블록에 들어가지 않고 기다리는 거래)의 거래, 오프체인(블록체인 밖의) 주문 장부, 중앙화 거래소 계정, 브리지는 자료 밖에 있어요. 한 운영자가 다루는 주소들을 한 주체로 합치려고 하지도 않아요. 다섯 토큰이 아닌 토큰은 빼요. 기본 이더(ether)는 WETH로만 들어와요.

\end{itemize}

\dissertationSubheading[sec:collection-strategy]{두 단계로 자료 모으기}

자료 모으기와 평가는 서로 떨어진 두 단계로 나누어 해요. 그래서 표, 벤치마크, 강건성 실행을 클라우드 비용(인터넷으로 빌려 쓰는 컴퓨터에 내는 돈) 없이 필요한 만큼 몇 번이든 다시 만들 수 있어요. 우리도 할 수 있고, 추출을 돌린 사람이면 누구나 할 수 있어요.

\begin{itemize}
\item 1단계(평판 이벤트 흐름): 스캔~A는 2023년 10월 1일부터 2026년 6월 30일까지의 허락, 송금, GMX 이벤트를 이 연구의 프로젝트로 복사해요. 스캔~B는 W1의 등록 문서가 커밋되면 2026년 7월부터 9월까지의 이벤트를 복사해요(\ref{sub:fresh-holdout}절). 질의 하나는 적어도 세 주인에게서 허락을 받은 스펜더를 목록으로 뽑아요. 코드 점검으로 계획의 코호트를 정하고, 또 다른 질의로 그 코호트의 에고 이벤트를 뽑아요. 주 분석은 이 이벤트들 가운데 다섯 토큰의 것만 남겨요. 그리고 각 동결일의 코호트를 그 점수 기간의 허락에서 뽑고, 그 코호트의 이벤트만 남겨요. 각 기준 시점, 곧 $T_{\mathrm{obs}}$와 $t_1$에서는 다음을 해요. 기준 시점의 점수 기간에 있는 허락을 (토큰, 주인, 스펜더) 세 개 묶음마다 (블록 번호, 로그 번호) 순서로 가장 최근 허락 하나로 줄이고, 그 허락이 0보다 크면 남겨요. 그 기간의 송금은 보낸 쪽--받는 쪽 쌍마다 모아요. 그리고 금액을 달러로 매겨요(\ref{sec:amount-revision}절). 대리 지표와 홀드아웃 라벨도 함께 계산해요. 이 표들이 그래프 만들기에 들어가는 유일한 입력이에요.

\item 2단계(로컬 평가, 연구자 자신의 컴퓨터에서 하는 평가): 로컬 컴퓨터는 내려받은 1단계 표로 GMX 닫기를 해독하고(풀어 읽고), 점수, 일치도 통계, 홀드아웃, 벤치마크, 강건성 실행을 계산해요. 이때 빅쿼리에는 질의하지 않아요.

\item 비용 안전장치: 모든 질의는 먼저 미리 돌려 보기(dry run, 실제로 돌리지 않고 비용만 어림해 보기)를 해요. 어림한 비용 때문에 쓴 돈이 이 연구의 예산 USD~$50$을 넘게 되는 질의는, 실행하는 사람이 추가 플래그(명령에 따로 붙이는 옵션)로 확인하지 않으면 중단돼요.\footnote{이 예산은 이 연구가 만든 처리 과정(파이프라인)이 스스로 지키게 해요. 빅쿼리가 걸어 두는 제한이 아니에요.} 모든 작업은 요금이 매겨진 바이트 수, 그리고 비용과 함께 비용 장부에 적혀요(부록~\ref{ch:appendix-eval}). 그래서 자료를 얻는 데 든 비용을 나중에 검사해 볼 수 있어요.

\end{itemize}

\dissertationSubheading[sec:pipeline-architecture]{처리 과정의 짜임}

그림~\ref{fig:pipeline-architecture}에 공개 로그에서 연구 표까지 이어지는 처리 과정 전체를 그렸어요. 원격 단계(클라우드에서 하는 단계)는 모두 이 연구의 빅쿼리 프로젝트 안에 표를 써요. 내려받은 파일은 한 번도 바꾸지 않아요.

\begin{figure}[htbp]
\centering
\begin{adjustbox}{max width=\linewidth}
\begin{tikzpicture}[
  font=\small,
  node distance=1.0cm and 1.0cm,
  box/.style={draw, rounded corners=2pt, align=center, minimum height=0.95cm, minimum width=2.2cm, inner sep=3pt},
  wide/.style={draw, rounded corners=2pt, align=center, minimum height=0.95cm, minimum width=2.6cm, inner sep=3pt},
  arrow/.style={-{Stealth[length=2.5mm]}, thick}
]
% Top row: remote acquisition (left to right)
\node[wide] (bq) at (0,0) {BigQuery\\공개 로그};
\node[wide] (extract) at (3.2,0) {로그 뽑기\\(스캔 A와 B; 토픽과\\코호트로 거르기)};

% Local evaluation. Row 1: events -> edge lists -> scores -> evaluation -> export.
% Row 2: proxy metrics from the same events. Row 3: temporal holdout.
\node[box] (decode) at (0,-2.3) {표 불러오기;\\GMX 해독};
\node[box] (prep) at (3.0,-2.3) {화살표 목록\\만들기};
\node[wide] (ranks) at (6.2,-2.3) {페이지랭크 점수\\(엔도스랭크, AWP,\\C-PR, S-PR)};
\node[wide] (align) at (9.5,-2.3) {일치도와\\벤치마크};
\node[wide] (tabexp) at (12.7,-2.3) {표와 그림\\내보내기};
\node[box] (proxy) at (3.0,-3.9) {대리\\지표};
\node[wide] (hold) at (6.2,-5.2) {시간 홀드아웃\\(W0과 W1)};

\draw[arrow] (bq) -- (extract);
\draw[arrow] (extract.south) -- ++(0,-0.55) -| (decode.north);
\draw[arrow] (decode) -- (prep);
\draw[arrow] (prep) -- (ranks);
\draw[arrow] (ranks) -- (align);
\draw[arrow] (align) -- (tabexp);
\draw[arrow] ([xshift=4pt]decode.south) |- (proxy.west);
\draw[arrow] ([xshift=-4pt]decode.south) |- (hold.west);
\draw[arrow] (proxy.east) -| (align.south);
\draw[arrow] (hold.east) -| (tabexp.south);

\node[draw, dashed, rounded corners, fit=(bq)(extract), inner sep=8pt, label=above:{\footnotesize 원격 / 자료 얻기}] {};
\node[draw, dashed, rounded corners, fit=(decode)(prep)(ranks)(align)(tabexp)(proxy)(hold), inner sep=8pt, label=below:{\footnotesize 로컬 / 재현 가능}] {};
\end{tikzpicture}
\end{adjustbox}
\caption[{처리 과정 한눈에 보기: 공개 로그에서 점수, 점검, 표까지.}]{처리 과정을 한눈에 본 그림이에요. 허락, 송금, GMX 로그는 원격에서 두 번의 스캔으로 복사하고, 코호트에 대한 쌍 표로 모아요. 로컬에서는 GMX 닫기를 해독하고, 쌍 표를 화살표 목록으로 바꿔요. 이 목록으로 엔도스랭크, AWP, C-PR, S-PR을 계산해요. 대리 지표는 같은 이벤트에서 나란한 다른 갈래로 나와요. 그 뒤에 일치도 계산과 벤치마크 계산, 그리고 시간 홀드아웃이 이어져요. 결과 표와 그림은 맨 마지막에 내보내요.}
\label{fig:pipeline-architecture}
\end{figure}

단위 테스트(코드의 작은 부분마다 맞게 도는지 보는 시험)는 점수를 매기는 코드와 통계 코드를 일부러 만든 시험용 자료 위에서 돌려요.

\dissertationSubheading[sec:event-decoding]{이벤트 해독}

스마트 계약이 남기는 로그에는 토픽과 데이터 내용이 있지만, 칸 이름은 없어요. 해독은 각 줄을 종류가 정해진 기록으로 바꿔요. 이 기록은 그래프를 만들고 점검을 하는 데 써요. 여기서 잘못이 생기면 두 그래프가 아무도 모르게 망가질 거예요. 토픽 해시와 칸 대응은 아래에 적었어요.

\begin{itemize}
\item 로그의 생김새: 각 로그에는 로그를 내보낸 계약, 토픽 여러 개, 그리고 ABI(계약의 명령과 이벤트를 적는 방식)로 적은 데이터 덩어리가 기록돼요. 첫 번째 토픽은 이벤트 서명(이벤트 이름과 값 종류를 적은 글)의 keccak-256 해시예요. 그 뒤로 색인된 값(찾기 쉽게 토픽에 넣는 값)이 $32$바이트 단어로 이어져요(주소는 왼쪽을 채워 길이를 맞춰요). 색인되지 않은 값은 ABI 순서대로 데이터 칸에 들어가요\parencite{vogelsteller2015erc20}. ERC-721 토큰도 같은 서명으로 \texttt{Transfer}와 \texttt{Approval}을 내보내요. 하지만 토큰 번호(ID)를 네 번째 토픽에 넣어요. 그래서 토픽이 정확히 세 개여야 한다는 조건이 이 토큰들을 걸러 내요. SQL은 토픽으로 거르고 ERC-20 칸을 해독해요. GMX 이벤트는 파이썬(Python)이 해독해요.

\item ERC-20 \texttt{Approval}: 이 이벤트의 토픽 해시는 아래와 같아요.
\begin{center}
\texttt{0x8c5be1e5ebec7d5bd14f71427d1e84f3dd0314c0f7b2291e5b200ac8c7c3b925}
\end{center}
이 값은 \texttt{Approval(address,address,uint256)}의 keccak-256 해시예요. 주인과 스펜더는 토픽~$1$과~$2$에 들어 있어요. 허락 금액은 데이터 칸에서 토큰의 원래 단위(소수점 조정을 하지 않은 가장 작은 단위)로 해독하고, 달러로 값을 매겨요(\ref{sec:amount-revision}절). 금액이 0인 허락도 남겨 둬요. 그 허락은 지금 효력이 있는 허락을 0으로 되돌리기 때문이에요. \texttt{permit}도 같은 \texttt{Approval} 이벤트를 내보내요. 그래서 허락은 온체인 \texttt{approve()} 호출로 생겼든, 오프체인 서명으로 생겼든 똑같이 세어요\parencite{lundfall2020permit}.

\item ERC-20 \texttt{Transfer}: 이 이벤트의 토픽 해시는 아래와 같아요.
\begin{center}
\texttt{0xddf252ad1be2c89b69c2b068fc378daa952ba7f163c4a11628f55a4df523b3ef}
\end{center}
이 값은 \texttt{Transfer(address,address,uint256)}의 keccak-256 해시예요. 보낸 쪽과 받는 쪽은 색인되어 있어요. 금액은 그날의 가격으로 달러 값을 매겨요. 금액이 0인 송금과, 토큰의 첫 가격이 나오기 전에 한 송금은 버려요. 발행 로그와 소각 로그(토큰을 새로 만들거나 없앨 때 남는 로그)는 영 주소(모든 자리가 0인 주소)를 써요. 이 로그들은 코호트 계약이 끼어 있으면 언제나 남겨요. 자기 송금(자기 자신에게 보낸 송금)은 이벤트 표에 남아 있어요. 하지만 들어오는 상대를 세는 계산에서는 빼요(\ref{sec:proxy-formulas}절). 자기 송금으로는 거래 상대가 드러나지 않기 때문이에요.

\item GMX~V2 포지션 닫기: 닫기 하나는 GMX~V2 이벤트 발신 계약(이벤트를 내보내는 계약)의 \texttt{PositionDecrease} 로그 하나예요. 그래서 일부만 닫는 것과 청산(손실이 너무 커져 거래가 강제로 끝나는 일)도 닫기로 세어요. 파이썬은 닫기마다 계정, 기본 손익(번 돈이나 잃은 돈), 주문 종류를 해독해요. 계약 계정의 닫기는 모두 오류 없이 해독되었어요.

\item 주소 맞추기와 단위: 주소는 표를 잇기(조인) 전에 소문자로 바꿔요. 금액은 \texttt{uint256} 값을 배정밀도 실수로 바꾼 것이에요. 소수 자릿수가 $q$인 토큰이면 이 값을 $10^{q}$로 나누고, 토큰의 달러 가격을 곱해요. 그래서 무제한 허락, 곧 원래 단위로 $2^{256}-1\approx1.16\times10^{77}$인 허락은 아주 큰 달러 금액이 돼요(그 효과는 \ref{sub:endorserank-graph}절과 \ref{sub:coupled-graph}절에 나와요).

\end{itemize}

\dissertationSubsubheading[sec:amount-revision]{토큰, 가격, 그리고 금액 수정}

분석 계획은 세 가지를 정해 두었어요. 금액은 원래 단위로 쓰고, 모든 ERC-20 토큰을 쓰고, 식~\eqref{eq:value-transform}에는 $b=1$을 쓰기로 했어요. 이 설정에는 문제가 세 가지 있어요. 첫째, 금액이 원래 단위로 $10$ 이상이면 $V(z)$와 $1$의 차이는 $10^{-4}$ 안이에요. 그리고 거의 모든 송금과 허락이 그 크기에 이르러요. 그래서 금액은 엔도스랭크에서도 AWP에서도 아무 역할을 하지 못했어요. 둘째, C-PR의 층들은 소수 자릿수가 서로 다른 토큰들의 원래 단위를 그대로 더했어요. 그리고 무제한 허락은 그런 허락을 가진 모든 주인의 줄을 거의 다 차지했어요. 셋째, 어떤 토큰 계약이든, 실제로는 아무 일에도 끼지 않은 주인이나 보낸 쪽을 적은 \texttt{Approval} 로그와 \texttt{Transfer} 로그를 내보낼 수 있어요. 그래서 아무 토큰이나 남긴 로그는 보증이나 지불의 증거가 되지 못해요.

그래서 수정 분석은 등록 분석의 결과를 안 뒤에 정의되었어요. 수정 분석은 금액을 달러로 매기고, 토큰 다섯 개만 남기고, 각 동결일 전 12개월의 자료로 각 계약에 점수를 매겨요(\ref{sec:data-sources}절). 토큰이 자격을 얻으려면, 아비트럼 원에 대한 체인링크(Chainlink)의 데이터 피드 목록에 그 토큰이나 그 토큰이 나타내는 자산의 시장 가격 피드(블록체인에 올라오는 가격 정보)가 시장 위험 등급 낮음이나 보통으로 올라 있어야 해요\parencite{chainlinknddatafeeds}. 자격을 얻은 토큰은 DefiLlama에 그 토큰의 하루 가격이 있으면 가격이 있는 토큰으로 쳐요\parencite{defillamanddocs}. 이 목록이 피드와 짝지어 둔 토큰 $133$개 가운데 $64$개가 자격을 얻었고, 그 가운데 $46$개에 가격이 있어요. 분석은 가격이 있는 토큰 가운데 $T_{\mathrm{obs}}$의 점수 기간에서 기록이 가장 많은 다섯 개를 남겨요. 바로 WETH, USDC, USDT0, WBTC, ARB예요. 송금은 그 송금이 있었던 UTC 날의 DefiLlama 하루 종가(그날 마지막 가격)로 값을 매겨요. 마지막 허락(가장 최근 허락)은 기준 시점이 있는 날의 가격으로 값을 매겨요. $V$의 기울기는 $b=\ln 3/m$이에요. 그래서 $m$달러의 금액은 무게가 절반이에요. $m$은 W0 점수 기간의 달러 금액 중앙값이에요. 그래서 $t_1$ 뒤의 어떤 이벤트도 이 값에 들어가지 않아요. 송금은 $m_T=170$달러이고, $t_1$에서 무제한이 아닌 마지막 허락은 $m_A=100$달러예요. C-PR의 허락 층에서 허락 하나는 많아야 $U=19{,}000$달러까지만 세어요. 이 값은 같은 기간에서 송금 하나하나의 금액이 갖는 99번째 백분위예요. 무제한이 아닌 허락 자체의 99번째 백분위는 약 $10^{63}$달러예요. 그래서 허락에서 상한을 잡으면, 여전히 무제한 허락이 줄을 좌우하게 돼요. 코호트는 이미 뽑아 둔 에고 이벤트에서 다섯 토큰으로 뽑아요. 새로 스캔하지 않아요.

규칙 A와 B, 그리고 두 종류의 화살표 어느 것으로도 만들지 않은 결과에 대한 비교는 그 기준값과 결론 문구 그대로 고친 점수에 적용해요. 이 점수들은 등록 분석의 결과를 안 뒤에 정의되었어요. 그래서 이 판정은 재분석이고, 등록한 시험이 아니에요. 부록~\ref{sec:map-registered}은 등록 분석을 수정 분석 옆에 나란히 놓아요. 부록~\ref{sec:data-tokens}에는 토큰, 빠진 기록, 가격 점검이 나와요.

\dissertationSubheading[sec:graph-construction]{그래프 만들기}

엔도스랭크와 AWP는 화살표에 무게를 다는 방법이 같아요. 또 같은 페이지랭크 풀이 프로그램, 곧 식~\eqref{eq:pagerank}을 돌려요. 다른 점은 두 가지예요. 화살표를 만드는 이벤트가 다르고, 걷기가 다시 출발하는 곳이 달라요. 각 방법은 양 끝 가운데 적어도 하나가 씨앗 집합 $\mathcal{W}$에 있는 화살표를 남겨요. 그래서 씨앗과 화살표 하나로 바로 이어진 이웃도 점으로 들어와요. 그리고 각 방법은 평가하는 모든 계약에 점수를 내놓아요. 부분 그래프에 없는 계약은 $0$점을 받아요.

엔도스랭크의 화살표는 보증을 나타내고, AWP의 화살표는 송금을 나타내요(\ref{sec:allowance-endorsement}절). 두 방법 모두 이벤트가 얼마나 오래되었는지에 따라 무게를 달아요. 그래서 최근 이벤트일수록 더 크게 세어져요. 또 두 방법 모두 이벤트의 달러 금액을 상한이 있는 함수에 넣어요. 그래서 금액이 클수록 더 크게 세어지지만, 이벤트 하나는 많아야 한 번만큼만 세어져요. 엔도스랭크가 AWP에서 무게를 가져오기 때문에, AWP를 먼저 설명해요. 두 종류의 화살표를 모두 걷는 C-PR은 맨 마지막에 설명해요.

\dissertationSubsubheading[sub:awp-graph]{AWP 송금 그래프}

AWP는 \textcite{do2023awp}의 방법이고, IEEE RIVF 2023 학회에서 발표되었어요. 송금 $e$의 금액이 $x_e$달러이고, 이 송금이 점수 기간이 끝나는 기준 시점($T_{\mathrm{obs}}$, W0에서는 $t_1$)보다 $\Delta t_e$일 전에 일어났다고 해요. 이 송금은 보낸 쪽에서 받는 쪽으로 가는 화살표에 $\sigma(\Delta t_e)\,V(x_e)$의 무게로 들어가요. 시간 무게는 아래의 로지스틱 감쇠예요.
\begin{equation}
\label{eq:logistic-decay}
\sigma(\Delta t)
=
\frac{1}{1+\exp\bigl(k(\Delta t-t_0)\bigr)},
\qquad
k=0.01,\quad t_0=180\text{일},
\end{equation}
금액 무게는 아래와 같아요.
\begin{equation}
\label{eq:value-transform}
V(z)
=
\frac{2}{1+e^{-bz}}-1,
\qquad
b=\frac{\ln 3}{m},
\end{equation}
이 값은 $z=0$일 때 $0$이고, 중앙값 금액 $z=m$일 때 $1/2$이에요. 그리고 금액이 커질수록 $1$에 가까워져요. 송금에서는 $b=b_T=\ln 3/m_T$예요(\ref{sec:amount-revision}절). 화살표 무게는 송금들의 무게를 더한 값이에요.
\begin{equation}
\label{eq:transfer-weight}
w_T(u,v)
=
\sum_{e:\,u\to v}
\sigma(\Delta t_e)\,V(x_e).
\end{equation}
이 합은 기간 안에 $u$에서 $v$로 간 모든 송금을 토큰과 상관없이 더해요. 자기 송금($u=v$)도 버리지 않아요. 이 송금은 $u$ 위의 자기 고리(자기에게 다시 돌아오는 화살표)가 돼요. 걷기는 주소의 활발함에 비례해서 그 주소에서 다시 출발해요.
\begin{equation}
\label{eq:activeness}
X_u
=
\sum_{v\neq u}\;
\max_{e:\,u\to v}\sigma(\Delta t_e),
\end{equation}
이 식은 $u$가 송금을 보낸 주소마다, 그 주소로 보낸 마지막 송금의 시간 무게를 골라 모두 더해요. 식~\eqref{eq:pagerank}의 재시작 분포는 $\pi_u=X_u/\sum_{v\in V_T}X_v$예요. 매달린 점에 고인 몫도 이 분포가 받아서 나눠요. 다른 주소에 한 번도 보내지 않은 주소는 $X_u=0$이에요. 그래서 그 주소는 다시 출발 몫을 받지 못하고, 점수는 들어오는 화살표에서만 나와요. 감쇠 계수는 $d=0.85$예요.

그 논문은 이더리움 거래에 점수를 매기고, 매개변수 값은 알려 주지 않아요. 그 논문의 시간 무게는 식~\eqref{eq:logistic-decay}의 로지스틱 함수와 같아요. 다만 나이 대신 달력 시간으로 적혀 있어요. $k$, $t_0$, $b$는 이 연구가 고른 값이에요. 중앙값 금액의 송금은 무게가 절반이고, 중앙값의 두 배인 송금은 무게가 $0.8$이에요. 그래서 화살표 무게는 송금 횟수뿐 아니라 보낸 금액에 따라서도 커져요. 점수 기간 안에서 로지스틱 감쇠는 오래된 이벤트를 크게 깎아요(그림~\ref{fig:logistic-decay-methods}). 무게는 $\sigma(0)=0.858$에서 $\sigma(182)=0.495$를 거쳐, 기간 첫날의 이벤트에서는 $\sigma(365)=0.136$까지 떨어져요.

\begin{figure}[htbp]
\centering
\includegraphics[width=0.8\textwidth]{\figpath{logistic-decay.pdf}}
\caption[{AWP와 엔도스랭크의 로지스틱 감쇠 $\sigma(\Delta t)$($k=0.01$, $t_0=180$일).}]{AWP와 엔도스랭크의 로지스틱 감쇠 $\sigma(\Delta t)$예요($k=0.01$, $t_0=180$일). 가로축은 이벤트의 나이, 곧 이벤트가 기준 시점보다 며칠 전에 일어났는지를 보여 줘요. 기준 시점은 $T_{\mathrm{obs}}=$ 2026년 6월 30일이고, W0에서는 $t_1=$ 2026년 3월 31일이에요. 세로축은 송금이나 허락의 시간 무게를 보여 줘요. 색칠한 띠는 두 동결일 어느 쪽에서든 나올 수 있는 나이(0일부터 $365$일까지)를 표시해요. 이 띠 안에서 무게는 $0.86$에서 $0.14$까지 떨어져요.}
\label{fig:logistic-decay-methods}
\end{figure}

알고리즘~\ref{alg:awp}에 그래프를 만들고 점수를 매기는 단계를 형식을 갖추어 적었어요.

\begin{algorithm}[htbp]
\caption{AWP 그래프 만들기와 점수 매기기}
\label{alg:awp}
\begin{algorithmic}[1]
\Require 달러 금액이 붙은 송금 $\mathcal{T}$, 씨앗 계약 $\mathcal{W}$, 관측 기간의 끝 $T_{\mathrm{obs}}$, 매개변수 $k,t_0,b_T,d,\varepsilon,I_{\max}$
\Ensure $\mathcal{W}$에 대한 점수 $\mathbf{s}_{\mathrm{AWP}}$
\State $E \gets \emptyset$; 모든 주소 $u$에 대해 $X_u\gets 0$
\For{$\mathcal{T}$ 안의 각 송금 $e=(u,v,x_e,t_e)$에 대해}
  \State $\Delta t \gets (T_{\mathrm{obs}}-t_e)$ (일 단위)
  \State 화살표 $(u\to v)$에 $\sigma(\Delta t)\,V(x_e)$를 더해 쌓기 \Comment{식~\eqref{eq:logistic-decay}--\eqref{eq:transfer-weight}; $u=v$이면 자기 고리가 됨}
\EndFor
\For{각 보낸 쪽 $u$에 대해}
  \State $X_u \gets \sum_{v\neq u}\max_{e:\,u\to v}\sigma(\Delta t_e)$ \Comment{식~\eqref{eq:activeness}}
\EndFor
\State 무게가 0보다 크지 않은 화살표 지우기; $\mathcal{W}$에 닿는 화살표만 남기기
\State 남은 화살표의 점 $u$마다 $\pi_u \gets X_u/\sum_{v}X_v$
\State $\mathbf{s} \gets \mathrm{WeightedPageRank}(E, \boldsymbol{\pi}, d, \varepsilon, I_{\max})$ \Comment{식~\eqref{eq:pagerank}--\eqref{eq:transition-impl}}
\State \Return $i\in\mathcal{W}$마다 $\mathbf{s}_{\mathrm{AWP}}[i] \gets \mathbf{s}[i]$ (그 밖에는 $0$)
\end{algorithmic}
\end{algorithm}

\dissertationSubsubheading[sub:endorserank-graph]{엔도스랭크 허락 그래프}

엔도스랭크는 AWP의 화살표 무게를 허락 그래프에 얹어요. 미리 손질(전처리)을 마친 마지막 허락들의 집합을 $\mathcal{A}$라고 해요. (토큰, 주인--스펜더) 세 개 묶음마다, 점수 기간 안의 가장 최근 \texttt{Approval} 이벤트 하나만 남겨요. 이때 순서는 (블록 번호, 로그 번호) 순서로 정해요. 허락 금액이 0이면 그 화살표를 지워요. 이렇게 남은 허락 $a^{\star}(u,v,\mathrm{token})$은 달러 금액이에요. 이 허락은 $T_{\mathrm{obs}}$보다 $\Delta t^{\star}(u,v,\mathrm{token})$일 전에 정해졌어요. 그리고 그 나이와 그 금액을 가진 이벤트 하나로 화살표에 들어가요.
\begin{equation}
\label{eq:approve-weight}
w_A(u,v)
=
\sum_{\mathrm{token}}
\sigma\bigl(\Delta t^{\star}(u,v,\mathrm{token})\bigr)\,
V\bigl(a^{\star}(u,v,\mathrm{token})\bigr),
\end{equation}
여기서 $\sigma$, $V$, $k$, $t_0$은 AWP와 같고, $b=b_A=\ln 3/m_A$예요. 방향이 있는 화살표 $u\to v$는 주인 $u$가 스펜더 $v$를 보증한다는 뜻이에요. 중앙값 금액의 허락은 시간 무게의 절반으로 들어가고, 큰 허락이나 무제한 허락은 시간 무게 전체에 가깝게 들어가요. 그래서 $w_A(u,v)$는 $u$가 $v$에게 허락한 금액이 클수록 커져요. 하지만 허락 하나가 1보다 크게 세어지지는 않아요. 감쇠 계수는 $d=0.85$예요.

걷기는 고르게 다시 출발해요. AWP의 다시 출발 규칙을 허락에 옮기면, 걷기는 주인 $u$에서 $X^A_u=\sum_{v\neq u}\max_{\mathrm{token}}\sigma\bigl(\Delta t^{\star}(u,v,\mathrm{token})\bigr)$에 비례해서 다시 출발해요. 이 합은 $u$의 0보다 큰 마지막 허락을 받은 스펜더 $v$들에 걸쳐 더해요. 그래서 이 무게는 $u$가 최근에 허락해 준 스펜더가 많을수록 커져요. 그러면 보증을 해 주는 쪽이 상을 받게 돼요. 그런데 엔도스랭크는 보증을 받는 계약에 순위를 매기려는 방법이에요. 분석 계획은 고르게 다시 출발하기로 정했어요. 그리고 이 변형, 곧 활동 재시작을 쓰는 엔도스랭크를 엔도스랭크 옆에 함께 보고할 것으로 정해 두었어요(표~\ref{tab:endorserank-restarts}).

이 무게는 마지막 허락의 금액과 시간을 담아요. 기간이 끝날 때 실제로 남아 있는 허락 금액은 아니에요. 한 쌍 안에서 토큰마다, 가장 새 허락이 그전 허락들을 대신해요\parencite{vogelsteller2015erc20}. 점수 기간이 시작되기 전에 마지막으로 정해진 허락은, 아직 효력이 있어도 화살표가 되지 않아요. 그래서 허락 그래프에는 동결일 전 1년 동안 주인들이 해 주었거나 새로 고친 허락만 담겨요. 주인이 유니스왑(Uniswap)의 Permit2 같은 허락 관리 계약(허락을 모아서 대신 관리해 주는 계약)에 허락하면, 그 주인은 그 계약으로 가는 화살표로만 나타나요. 나중에 그 계약을 거쳐 토큰을 쓰는 프로그램으로 가는 화살표는 나타나지 않아요. 어떤 토큰은 \texttt{transferFrom} 안에서도 줄어든 허락 금액으로 \texttt{Approval}을 내보내요. 이런 토큰에서는 마지막 기록이 스펜더가 마지막으로 쓰고 남은 금액일 수 있고, 그 날짜는 그렇게 쓴 날이에요. 그래서 시간 무게가 쓴 일을 일부 기록하게 돼요. 또 다 써 버린 허락은 0으로 끝나서 빠져요. 그리고 점수 기간 전에 해 주고 라벨 기간에 쓴 허락이 새 허락 쌍(새로 생긴 주인–스펜더 허락)으로 나타날 수 있어요. 다른 토큰에서는 줄어든 허락도 처음 허락한 금액을 그대로 가져요. 마지막 기록이 이렇게 남은 금액인 경우가 얼마나 잦은지는 재지 않았어요.

그래프를 만들고 페이지랭크를 돌리는 단계는 알고리즘~\ref{alg:endorserank}에 적어 두었어요.

\begin{algorithm}[htbp]
\caption{엔도스랭크 그래프 만들기와 점수 매기기}
\label{alg:endorserank}
\begin{algorithmic}[1]
\Require 달러 금액이 붙은 마지막 허락 $\mathcal{A}$, 씨앗 계약 $\mathcal{W}$, 관측 기간의 끝 $T_{\mathrm{obs}}$, 매개변수 $k,t_0,b_A,d,\varepsilon,I_{\max}$
\Ensure $\mathcal{W}$에 대한 점수 $\mathbf{s}_{\mathrm{ER}}$
\State $E \gets \emptyset$
\For{$\mathcal{A}$ 안의 각 마지막 허락 $(u,v,\mathrm{token},a^{\star},t^{\star})$에 대해}
  \State $\Delta t^{\star} \gets (T_{\mathrm{obs}}-t^{\star})$ (일 단위)
  \State 화살표 $(u\to v)$에 $\sigma(\Delta t^{\star})\,V(a^{\star})$를 더해 쌓기 \Comment{식~\eqref{eq:approve-weight}}
\EndFor
\State 무게가 0보다 크지 않은 화살표 지우기; $\mathcal{W}$에 닿는 화살표만 남기기
\State $\mathbf{s} \gets \mathrm{WeightedPageRank}(E, \text{균등 }\boldsymbol{\pi}, d, \varepsilon, I_{\max})$ \Comment{식~\eqref{eq:pagerank}--\eqref{eq:transition-impl}}
\State \Return $i\in\mathcal{W}$마다 $\mathbf{s}_{\mathrm{ER}}[i] \gets \mathbf{s}[i]$ (그 밖에는 $0$)
\end{algorithmic}
\end{algorithm}

허락은 송금보다 드물어요. 그래서 매칭 코호트에서는 허락 부분 그래프가 더 듬성듬성해요(\ref{ch:results}장). 페이지랭크 반복 한 번에 드는 비용은 $O(|E|)$예요. 그래서 이 희소성(화살표가 적음)이 H3을 받치는 그래프 구조의 근거예요.

두 방법 모두 화살표 목록을 (출발, 도착, 무게) 세 개 묶음으로 저장해요. 겹치는 화살표는 무게를 더해서 하나로 합쳐요. 무게가 0보다 크지 않은 화살표는 풀기 전에 버려요.

\dissertationSubsubheading[sub:coupled-graph]{두 층을 함께 걷는 결합 연산자(C-PR)와 씨앗을 쓰는 빼 보기 실험(S-PR)}

엔도스랭크와 AWP는 각각 한 종류의 화살표 위에서만 걸어요. 그래서 각 방법은 계약을 한 가지 행동으로만 줄 세워요. 그 행동은 계약이 받은 쓸 권한이거나, 계약이 받은 지불이에요. 이 논문의 방법인 C-PR은 같은 주소들의 두 종류 화살표를 하나의 걷기 안에서 함께 걸어요. RQ4는 C-PR이 두 층 가운데 어느 쪽이든 한 층만 걷는 걷기보다 계약의 순위를 더 잘 매기는지 물어요. RQ5는 C-PR이 송금만 걷는 걷기보다 나아지는지 물어요. 분석 계획은 결합 연산자를, 각 이벤트에 금액으로 무게를 다는 두 층 위에서 정해 두었어요. 수정 분석에서는 이 금액을 달러로 매겨요(\ref{sec:amount-revision}절).
\begin{equation}
\label{eq:raw-layers}
\tilde w_A(u,v)
=
\sum_{\mathrm{token}}
\min\bigl(a^{\star}(u,v,\mathrm{token}),U\bigr),
\qquad
\tilde w_T(u,v)
=
\sum_{e:\,u\to v}
x_e\,\sigma(\Delta t_e).
\end{equation}
허락 층에는 엔도스랭크의 화살표가 있고, 송금 층에는 AWP의 화살표가 있어요. 다만 금액 무게가 컴퓨터의 소수 계산(부동소수점)에서 0이 되는 쌍은 예외예요(\ref{sub:cohorts}절). 이런 쌍은 C-PR의 층에만 남아 있어요. 그 밖에는 무게만 달라요. 금액은 상한이 있는 값 변환을 거치지 않고 들어가요. 허락 하나는 많아야 $U$까지만 세어요. 그리고 허락 층에는 로지스틱 감쇠가 없어요. C-PR을 위한 두 규칙은 모두 원래 단위로 된 이 층들 위에서 정했어요. 두 층은 합집합 점 집합 $V=V_A\cup V_T$ 위에서 그래프 하나를 이뤄요. 이동 행렬은 식~\eqref{eq:transition-impl}에서 나온 $P_A$와 $P_T$예요. 점 $u$의 결합된 줄은 두 층의 줄을 볼록 혼합(더해서 1이 되는 비율로 섞음)한 것이에요. 이때 비율은 $u$가 나가는 화살표를 가진 층들만 가지고 다시 맞춰요.
\begin{align}
P^{(\lambda)}_{u,\cdot}
&=
\alpha_A(u)\,P_{A,u,\cdot}+\alpha_T(u)\,P_{T,u,\cdot},
\label{eq:coupled-transition}\\
\alpha_A(u)
&=
\frac{\lambda\,\mathbf{1}[o_A(u)>0]}{D_{\lambda}(u)},
\qquad
\alpha_T(u)=\frac{(1-\lambda)\,\mathbf{1}[o_T(u)>0]}{D_{\lambda}(u)},
\nonumber
\end{align}
여기서 $D_{\lambda}(u)=\lambda\,\mathbf{1}[o_A(u)>0]+(1-\lambda)\,\mathbf{1}[o_T(u)>0]$이에요. $o_A(u)$와 $o_T(u)$는 두 층에서 $u$의 나가는 무게의 합이에요. 한 층에만 나가는 화살표가 있는 점은 무게 1로 그 층을 따라가요. 무게가 $0$인 층은 그 화살표와 점까지 함께 빠져요. $D_{\lambda}(u)=0$인 점은 매달린 점이고, 그 몫은 식~\eqref{eq:pagerank}의 세 번째 항이 맡아요. 순간 이동은 고르게 하고, $d=0.85$예요. $\mathbf{s}_{\mathrm{CPR}}$은 $P=P^{(\lambda)}$로 둔 식~\eqref{eq:pagerank}의 페이지랭크 벡터예요. 풀이 설정은 한 층 방법들과 같아요. 그림~\ref{fig:cpr-schematic}는 이 걷기의 한 걸음을 보여 줘요. 알고리즘~\ref{alg:cpr}에는 그래프를 만들고 점수를 매기는 단계를 적었어요.

\begin{figure}[htbp]
\centering
\begin{tikzpicture}[
  font=\small,
  vertex/.style={draw, circle, minimum size=0.9cm},
  allow/.style={-{Stealth[length=2.5mm]}, thick},
  xfer/.style={-{Stealth[length=2.5mm]}, thick, dashed}
]
\node[vertex] (x) at (0,1.5) {$x$};
\node[vertex] (u) at (0,0) {$u$};
\node[vertex] (y) at (0,-1.5) {$y$};
\node[vertex] (v) at (5.2,0.85) {$v$};
\node[vertex] (w) at (5.2,-0.85) {$w$};
\draw[allow] (x) -- node[above, font=\footnotesize] {$1$} (v);
\draw[allow] (u) -- node[above, font=\footnotesize, pos=0.42] {$\lambda$} (v);
\draw[xfer] (u) -- node[below, font=\footnotesize, pos=0.42] {$1-\lambda$} (w);
\draw[xfer] (y) -- node[below, font=\footnotesize] {$1$} (w);
\end{tikzpicture}
\caption[{결합 걷기 C-PR의 한 걸음.}]{결합 걷기 C-PR의 한 걸음이에요(식~\eqref{eq:coupled-transition}). 실선 화살표는 허락 화살표이고, 점선 화살표는 송금 화살표예요. 화살표에 붙은 값은 감쇠 계수를 곱하기 전의 이동 확률이에요. 주소 $u$는 두 층 모두에 나가는 화살표가 있어서, 두 층을 $\lambda$와 $1-\lambda$의 무게로 섞어요. $x$와 $y$는 한 층에만 나가는 화살표가 있어서, $\lambda$가 얼마이든 그 층을 따라가요. 걷기는 $1-d$의 확률로, 고르게 고른 주소에서 다시 출발해요.}
\label{fig:cpr-schematic}
\end{figure}

\begin{algorithm}[htbp]
\caption{C-PR 그래프 만들기와 점수 매기기}
\label{alg:cpr}
\begin{algorithmic}[1]
\Require 달러 금액이 붙은 마지막 허락 $\mathcal{A}$와 송금 $\mathcal{T}$, 씨앗 계약 $\mathcal{W}$, 관측 기간의 끝 $T_{\mathrm{obs}}$, 상한 $U$, 무게 $\lambda$, 매개변수 $k,t_0,d,\varepsilon,I_{\max}$
\Ensure $\mathcal{W}$에 대한 점수 $\mathbf{s}_{\mathrm{CPR}}$
\State 허락 층: $\mathcal{A}$ 안의 각 마지막 허락마다 $(u\to v)$에 $\min(a^{\star},U)$를 더해 쌓기 \Comment{식~\eqref{eq:raw-layers}}
\State 송금 층: $\mathcal{T}$ 안의 각 송금마다 $(u\to v)$에 $x_e\,\sigma(T_{\mathrm{obs}}-t_e)$를 더해 쌓기
\State 무게가 0보다 크지 않은 화살표 지우기; $\mathcal{W}$에 닿는 화살표만 남기기; 무게가 $0$인 층 빼기
\State $V \gets V_A\cup V_T$; $P_A, P_T \gets$ 줄마다 합이 1이 되게 맞춘 두 층 \Comment{식~\eqref{eq:transition-impl}}
\For{$D_{\lambda}(u)>0$인 각 점 $u\in V$에 대해} \Comment{나머지는 매달린 점}
  \State $P^{(\lambda)}_{u,\cdot} \gets \alpha_A(u)\,P_{A,u,\cdot}+\alpha_T(u)\,P_{T,u,\cdot}$ \Comment{식~\eqref{eq:coupled-transition}}
\EndFor
\State $\mathbf{s} \gets \mathrm{WeightedPageRank}(P^{(\lambda)}, \text{균등 }\boldsymbol{\pi}, d, \varepsilon, I_{\max})$ \Comment{식~\eqref{eq:pagerank}}
\State \Return $i\in\mathcal{W}$마다 $\mathbf{s}_{\mathrm{CPR}}[i] \gets \mathbf{s}[i]$ (그 밖에는 $0$)
\end{algorithmic}
\end{algorithm}

$\lambda=1$이면 이 연산자는 허락 층만 걷고, $\lambda=0$이면 송금 층만 걸어요. 이 두 걷기를 C-PR ($\lambda=1$)과 C-PR ($\lambda=0$)이라고 써요. 이 둘은 C-PR을 위한 두 규칙 모두에서 비교 대상이에요(RQ4와 RQ5; \ref{sec:holdout-protocol}절과 \ref{sub:fresh-holdout}절). 단위 테스트는 시험용 자료 하나에서, 양 끝 값이 각각 그 층만 고르게 다시 출발하며 걸었을 때의 점수를 내는지 점검해요. 이 점수 묶음은 양 끝에서 끊어져 있어요(연속이 아니에요). $0<\lambda<1$이면 점 집합은 $V_A\cup V_T$예요. 그리고 한 층에만 나가는 화살표가 있는 점은 $\lambda$가 얼마이든 그 층을 따라가요. 그런데 $\lambda=1$이면, 나가는 화살표가 모두 송금인 점은 매달린 점이 되거나 아예 없어져요. 그래서 $\lambda$가 $1$에 가까워질 때 C-PR은 그런 점의 송금 화살표를 여전히 따라가는 걷기에 가까워져요. 다른 쪽 끝에서도 마찬가지예요. $\lambda=0.5$(주 값)와 $\lambda\in\{0.25,0.75\}$(민감도 분석 값)는 계획의 설정 파일에서 정해 두었어요.

금액에 상한을 두지 않으면 허락 층이 아주 고르지 않게 될 거예요. $T_{\mathrm{obs}}$에서 0보다 큰 마지막 허락의 절반 가까이가 무제한 허락이에요(부록~\ref{sec:data-tokens}). 상한이 없으면 무제한 허락 하나하나가 그 주인의 줄을 거의 다 차지할 거예요. 그러면 두 종류의 허락을 함께 가진 주인의 무제한이 아닌 허락은 거의 아무 효과도 내지 못할 거예요. 상한이 있으면 무제한 허락은, 금액이 $U$달러 이상인 무제한이 아닌 허락과 무게가 같아요. 그리고 더 작은 허락들은 금액에 비례해 그 줄에서 제 몫을 지켜요. 엔도스랭크와 AWP는 이렇게 하는 대신 $V$로 모든 금액에 상한을 둬요.

S-PR 빼 보기 실험에서는 송금 층을 걸어요. 하지만 식~\eqref{eq:pagerank}의 고른 순간 이동 벡터 $\boldsymbol{\pi}$를 C-PR ($\lambda=1$)의 점수로 바꿔요. 이 점수는 $V_T$에 있는 점만 남기고, 그 안에서 합이 1이 되게 다시 맞춰요. $\mathbf{r}_{A}$를 $V_A$ 위에서 허락 층만 걸었을 때의 페이지랭크 벡터라고 해요. $u\notin V_A$인 점에는 $r_{A,u}=0$을 붙여 넓혀요. 그러면 다음과 같아요.
\[
\pi_u=\frac{r_{A,u}}{\sum_{v\in V_T} r_{A,v}}\qquad(u\in V_T),
\]
그래서 허락 점수가 없는 송금 그래프의 점은 다시 출발 몫을 하나도 받지 못해요. $V_T$의 점이 하나도 $V_A$에 들어 있지 않으면, 분모가 0이 돼요. 그때 $\boldsymbol{\pi}$는 $V_T$ 위의 고른 분포로 돌아가요. 같은 벡터가 매달린 점의 몫도 다시 나눠요. 이 결과는 트러스트랭크(TrustRank, 믿을 만한 씨앗에서 다시 출발하는 페이지랭크)를 만드는 방식과 같아요\parencite{gyongyi2004trustrank}. 다만 사람이 손으로 라벨을 붙인 씨앗 대신, 허락에서 얻은 씨앗을 써요. 허락은 걷기가 어디서 다시 출발할지 정하고, 송금은 걷기가 어떻게 퍼질지 정해요. S-PR과 C-PR을 견주면, 허락 화살표가 퍼뜨리는 길로서 보태는 것과 씨앗으로서 보태는 것을 나누어 볼 수 있어요. 하지만 S-PR도 AWP나 C-PR ($\lambda=0$)처럼, 송금 화살표가 없는 계약에는 모두 0점을 줘요. 이런 계약은 동결일마다 스펜더 여섯 개 가운데 한 개쯤이에요(\ref{sec:holdout-results}절). 허락 화살표는 이 스펜더들도 줄 세워요. 그래서 송금만 걷는 걷기를 이 스펜더들을 줄 세우는 점수와 견주는 대비에는 언제나 순위 매기기와 덮는 범위(커버리지)가 섞여 들어가요. 규칙~B의 F1도 마찬가지예요.

결합 연산자에 드는 비용은 화살표 집합이 정해요. C-PR은 $E_A\cup E_T$ 위를 걸어요. 그래서 C-PR 반복 한 번은 두 층에서 반복을 한 번씩 하는 것만큼 비용이 들어요. 실행 시간도 AWP와 같거나 더 길 것으로 예상해요. 그래서 주장~C5의 속도 결과는 엔도스랭크만의 것이에요. S-PR은 층마다 한 번씩 풀어야 해요. 두 연산자 모두 \ref{sec:computational-benchmark}절의 벤치마크 절차에 따라 시간을 재요.

\dissertationSubheading[sec:holdout-protocol]{시간 홀드아웃 절차}

홀드아웃 W0에서 점수는 W0 점수 기간, 곧 2025년 4월 1일부터 $t_1=$ 2026년 3월 31일 23:59:59~UTC까지의 이벤트를 써요. 라벨은 2026년 4월 1일 00:00:00~UTC부터 $t_2=$ 2026년 6월 30일 23:59:59~UTC까지의 이벤트로 세요. 코호트는 $n=7{,}914$개의 스펜더 계약이에요. 이 계약들은 $t_1$에, 다섯 토큰 가운데 하나에 대해 기간 안에 정해진 0보다 큰 마지막 허락을 적어도 하나 가지고 있어요. 이 자료는 W0 코호트의 에고 네트워크예요. W0 코호트는 \ref{sec:data-sources}절의 규칙으로 같은 기간의 허락에서 뽑아요. 그래서 $t_1$ 뒤의 어떤 이벤트도, W0 자료를 모을 때 중심이 되는 계약을 고르는 데 끼지 않아요. 홀드아웃 표의 모든 점수는 이 이벤트들만으로 계산해요. 여기에는 엔도스랭크, AWP, 세 가지 $\lambda$ 값과 두 끝 값에서의 C-PR, S-PR, 그리고 원시 차수 기준선 두 개가 들어가요. 로지스틱 감쇠의 기준 시점도 $t_1$이에요. 식~\eqref{eq:logistic-decay}의 $\Delta t$는 $t_1$에서 거꾸로 재요. 그래서 나이는 $365$일까지 가고, 무게는 $0.858$에서 $0.136$까지 내려가요.

이 연구에서는 스펜더마다 라벨을 두 개 세요. 구성개념마다 하나씩이에요. 그래서 미래 기간에서 한 층 방법은 각자 자기 구성개념으로 시험을 받고, 혼합 점수들은 두 구성개념 모두로 시험을 받아요.
\begin{enumerate}
\item 새 허락 쌍: 주인--스펜더 쌍 가운데, 라벨 기간에 0보다 큰 허락을 받았지만 $t_1$의 마지막 허락 스냅숏(그 시점에 0보다 큰 마지막 허락을 찍어 둔 목록)에는 없던 쌍의 수예요. 이 수는 스펜더의 것으로 쳐요. 이 라벨은 보증 구성개념을 재요.
\item 새 송금자(새로 송금을 보내온 주소): 라벨 기간 동안 스펜더에게 송금을 보낸 서로 다른 주소의 수예요. 단, 그 주소는 점수 기간이 시작한 날(2025년 4월 1일)부터 $t_1$ 사이에는 그 스펜더에게 송금을 보낸 적이 없어야 해요. 자기 송금은 빼요. 이 라벨은 흐름 구성개념을 재고, 첫 번째 라벨과 짝을 이루는 송금 쪽 라벨이에요.
\end{enumerate}
기준선으로는 각 스펜더의 받은 허락 수(0보다 큰 허락을 가진 서로 다른 주인의 수)와 송금 입차수(스펜더 자신을 뺀, 송금을 보내온 서로 다른 주소의 수)를 써요. 둘 다 $t_1$ 시점의 값이에요. 한 층 페이지랭크는 각각 이 지역 통계 가운데 하나를 매끄럽게 다듬어요. 표에 이 기준선들의 줄도 나오기 때문에, 각 페이지랭크가 자기 원시 차수 위에 얼마나 더 보태는지(늘어난 몫)를 바로 읽을 수 있어요. RQ3에 답하는 것은 바로 이 늘어난 몫이에요. 페이지랭크 계수 하나만으로는 RQ3에 답할 수 없어요. C-PR은 두 통계를 모두 매끄럽게 다듬어요. 그래서 C-PR은 두 기준선 모두와 견주어 읽어요.

불확실성은 \ref{sec:statistical-tests}절의, 계약을 다시 뽑는 짝지은 부트스트랩으로 구해요($B=400$, 모두가 함께 쓰는 번호 행렬 하나). 그래서 같은 라벨에서 두 계수 사이의 차이는 모두 짝지은 차이예요. 계획에서 정한 대비는 바로 아래에 나오는 규칙~A의 대비 두 개예요. 이 기간의 다른 대비는 모두 라벨을 본 뒤에 계산했어요. 각 페이지랭크와 그 자신의 차수를 견준 대비와 S-PR의 대비도 여기에 들어요. C-PR과 다른 모든 점수 사이의 대비는 표~\ref{tab:cpr-contrasts}에 있어요. 엔도스랭크와 AWP 사이의 차이는 부록~\ref{ch:appendix-er-awp}에 있어요.

분석 계획에서 정한 규칙~A는 C-PR을 자기의 두 층을 각각 한 층만 걷는 것과 견주어 판정해요. 이 규칙은 두 부분으로 되어 있어요. 주 부분은 W0을 써요. $\lambda=0.5$의 C-PR은 새 허락 쌍에서 C-PR ($\lambda=1$)보다 나쁘지 않아야 해요($\Delta\tau=\tau_{\mathrm{CPR}}-\tau_{\lambda=1}$의 구간에 0이 들어 있거나, 구간이 0보다 위에 있어야 해요). 또 새 송금자에서는 C-PR ($\lambda=0$)보다 나아야 해요($\tau_{\mathrm{CPR}}-\tau_{\lambda=0}$의 구간이 0보다 위에 있어야 해요). 보조 부분은 같은 기간 평가와 매칭 코호트를 써요. 조건은 두 가지예요. 송금 지표 가족(비슷한 지표 묶음)과 허락 지표 가족에서, C-PR의 점 추정값(구간이 아니라 값 하나로 낸 추정)이 두 층 가운데 더 나은 쪽의 95\% 구간 안에 들어야 해요. 그리고 시빌 안정성 지표 가족에서는 C-PR이 두 층을 모두 넘어야 해요. 표~\ref{tab:tau-diff-hybrid}의 짝지은 대비는 더 엄격한 점검으로 이 규칙 옆에 함께 보고해요. 하지만 이 대비가 규칙의 판정을 정하지는 않아요. 주 부분이 통과하지 못하면, RQ4, 곧 C-PR이 두 층 가운데 어느 쪽이든 한 층만 걷는 걷기보다 나은지에 대한 답은 “아니요”가 돼요. 그리고 그 결과는 구간과 함께 보고해요. RQ5, 곧 C-PR이 송금만 걷는 걷기보다 나아지는지는 규칙~A가 다루지 않아요. 그것은 규칙~B가 W1에서 시험해요(\ref{sub:fresh-holdout}절). 라벨을 한 번 본 뒤에는 $\lambda$를 조정하지 않아요. 규칙~A를 정했을 때 이 연구는 로그를 하나도 뽑지 않은 상태였어요. 고친 점수에 대한 이 규칙의 판정은 재분석이에요(\ref{sec:amount-revision}절).

처리 과정은 라벨 기간 동안 두 가지를 더 세요. 하나는 허락 취소, 곧 $t_1$에 허락이 있던 쌍 가운데 금액이 0인 허락을 받은 쌍이에요(어떤 토큰에서는 다 써 버린 허락이에요, \ref{sub:endorserank-graph}절). 다른 하나는 허락 뒤에 토큰을 밖으로 빼 가는 일(drain)을 찾는 어림 규칙(휴리스틱)이에요. 어느 쪽도 결론을 받치는 증거로 내놓지 않아요(\ref{sec:holdout-results}절).

\dissertationSubsubheading[sub:trader-holdout]{W0의 트레이더 결과}

스펜더 코호트는 거래하는 계약에 대해서는 알려 주는 것이 별로 없어요. 탐색적(미리 정하지 않고 살펴본) 실행 하나에서는 동결일 $t_1$을 $87$개의 GMX~V2 계정에 적용했어요. 이 계정들은 계약이고, 동결일 그래프에 나타나고, 2026년 4월부터 6월까지 포지션을 적어도 세 번 닫았어요. 라벨로는 그 기간의 성공 라벨을 썼어요. 곧 이익을 낸 닫기의 수, 실현 이익(닫기들의 0보다 큰 기본 이익을 더한 값이고, 손실은 빼지 않아요), 그리고 닫기 가운데 이익을 낸 비율과 손실이 아니었던 비율이에요. 청산된 닫기는 이익을 내지 않은 것으로 세요. 또 처리 과정이 청산을 손실로 세지 않으므로, 손실이 아닌 것으로도 세요. 스펜더 홀드아웃의 모든 점수를 미리 정한 규칙 없이 이 라벨들과 견주었고, 짝지은 다시 뽑기는 400번이었어요. \ref{ch:results}장은 이 실행을 탐색적인 것으로 보고해요. 이 실행에는 청산 라벨이 없어요. 그리고 계약 트레이더(GMX V2에서 거래하는 계약)가 아주 적어서, 구간이 모두 넓어요.

\dissertationSubsubheading[sub:fresh-holdout]{등록 기간 W1: 2026년 7월부터 9월까지}

W0 홀드아웃은 어떤 규칙도 묻지 않았던 것도 하나 보여 주었어요. C-PR을 그 송금 층만 걷는 걷기와 견주면, C-PR은 나중에 새 허락 쌍을 얻은 스펜더를 더 높게 줄 세웠어요. 그 차이의 구간은 0보다 위에 있었어요. 반면에 새 송금자에서는 C-PR이 더 나빴어요(등록 분석에서; 부록~\ref{sec:map-registered}). 우리는 라벨이 다 나온 뒤에야 이것을 알아챘어요. 그래서 발견이 아니라 가설로 치고, 시험을 정할 때 이 연구가 아직 뽑지 않았던 기간에서 다시 시험해요. W0 결과를 안 뒤에 쓴 등록 문서, 그리고 이 기간의 상태를 “등록됨”으로 적은 설정 파일은 2026년 10월 9일 16:49~UTC에 공개 저장소에 커밋했어요(커밋 \texttt{0d3a057}). 그리고 2026년 6월 30일 뒤의 로그를 읽는 유일한 질의인 스캔~B는 이 커밋 뒤에 돌았어요. 추출 스크립트는 두 파일이 로컬 변경 없이 커밋되어 있고 상태가 “등록됨”일 때만 돌아가요. 그렇지 않으면 돌기를 거부해요. 그리고 이 스크립트는 커밋과 두 파일의 해시를 추출 기록 파일(manifest)에 적어요(부록~\ref{ch:appendix-eval}).

방법은 바뀌지 않았어요. C-PR은 계획에서 정한 대로 $\lambda=0.5$에서 쓰고, 민감도 분석 값 두 개와 S-PR도 함께 써요. 풀이 프로그램 설정, SQL, 매칭 코호트의 추출 조건, GMX 해독 프로그램, 짝지은 부트스트랩(다시 뽑기 400번, 시드 42)은 W0에서 쓴 것과 같아요. 점수는 $T_{\mathrm{obs}}$, 곧 2026년 6월 30일 23:59:59~UTC에 동결해요. 그래서 이 점수는 같은 기간 평가의 점수예요. 등록 분석에 속하는 등록 문서에서는 이 점수가 2023년 10월 1일부터 뽑은 모든 이벤트를 써요. 수정 분석에서는 $T_{\mathrm{obs}}$의 점수 기간에 있는 이벤트를 써요(\ref{sec:data-sources}절). 라벨은 2026년 7월 1일부터 9월 30일까지 세요.

두 코호트를 평가해요. 크기는 \ref{sub:cohorts}절에 있어요. 스펜더 코호트는 W0에서처럼 정해요. 곧, 동결일에 0보다 큰 마지막 허락을 적어도 하나 가진 모든 스펜더 계약이에요. 이 코호트에는 W0 스펜더의 여섯 개 가운데 다섯 개쯤이 들어 있어요. 그래서 W1은 대체로 같은 계약들에 대한 더 나중의 라벨 기간이에요. 트레이더 코호트는 GMX~V2 계정 가운데 계약이고, 동결일 그래프의 어느 한 층에서라도 점이고, 라벨 기간에 포지션을 적어도 세 번 닫은 계정이에요. 스펜더는 두 라벨을 그대로 가져요. 트레이더는 청산 없이 닫은 비율을 라벨로 받아요. 이 값은 한 계약이 7월부터 9월까지 닫은 포지션 가운데 청산으로 끝난 비율을 1에서 뺀 값이에요. 이 자료 안에서 부도에 가장 가까운 것이 이 값이에요. 또 이 값은 두 그래프 어느 쪽으로도 만들지 않았어요. 그리고 비율이기 때문에, 거래를 더 많이 했다는 것만으로 계약에 상을 주지 않아요. 이 라벨을 고를 때는 이 연구의 어떤 점수도 청산 라벨과 견주어 본 적이 없었어요. 이 연구는 해독 단계에서 계약 계정의 닫기 가운데 청산이 얼마나 흔한지만 알고 있었어요(\ref{sub:cohorts}절).

규칙~B는 C-PR을 그 송금 층만 걷는 걷기인 C-PR ($\lambda=0$)과 견줘요. 이때 짝지은 차이 $\Delta\tau=\tau_{\mathrm{CPR}}-\tau_{\lambda=0}$를 써요.
\begin{itemize}
\item F1(스펜더, 새 허락 쌍): $\Delta\tau$의 95\% 구간이 0보다 위에 있어야 해요.
\item F2(스펜더, 새 송금자): 구간의 아래 끝이 $-\delta$보다 위에 있어야 해요.
\item F3(트레이더, 청산 없이 닫은 비율): 구간의 아래 끝이 $-\delta$보다 위에 있어야 해요.
\end{itemize}
세 가지가 모두 맞을 때만 C-PR이 송금만 걷는 걷기보다 나아졌다고 쳐요. 그래서 여러 번 시험한 것을 바로잡는 보정은 필요 없어요. 허용 폭 $\delta=0.01$은 등록 분석에서 W0의 F2 대비가 가진 짝지은 구간의 절반 너비를 소수 둘째 자리로 올림한 값이에요. 등록 문서가 이렇게 정했어요(부록~\ref{sec:map-registered}). 등록 분석에서는 스펜더 코호트가 허용 폭보다 작은 차이도 가려낼 수 있었어요. 그래서 허용 폭은 그 자료에 필요한 것보다 넓었어요. 수정 분석의 더 작은 코호트에서는 F2의 짝지은 구간 너비가 허용 폭과 거의 같은 크기예요(W1에서 절반 너비 $0.012$). 그리고 트레이더에서는 두 걷기가 트레이더를 거의 똑같이 줄 세울 때만 구간의 아래 끝이 $-\delta$를 넘을 수 있어요. 그래서 F3은 통과할 힘(검정력)이 거의 없었어요. F3의 약점 하나도 규칙과 함께 적어 두었어요. 두 점수 모두 청산을 예측하지 못하면, 두 계수가 모두 0 가까이에 있어도 F3이 통과할 수 있어요. 그래서 두 계수를 판정 옆에 함께 보고해요. 결과 세 개를 더 정해진 순서로 보고하지만, 이것들은 결정에 들어가지 않아요. 청산 없이 닫은 비율에서 C-PR과 C-PR ($\lambda=0$)을 견준 우월성 검사(더 나은지 보는 검사, F4), 같은 라벨에서 C-PR과 C-PR ($\lambda=1$)을 견준 것(F5), 그리고 새 기간에서 다시 돌린 규칙~A의 주 부분(F6)이에요.

민감도 분석 두 개도 같은 커밋에 등록했어요. 이 연구가 7월부터 9월까지의 로그를 하나도 뽑기 전이었어요. 둘 다 결정 옆에 함께 보고하고, 어느 쪽도 결정에 들어가지 않아요. 첫 번째는 C-PR ($\lambda=0$) 대신 AWP(\ref{sub:awp-graph}절)를 넣어 F1--F4를 되풀이해요. 두 번째는 그래프에서 빠진 코호트 계약에 0점을 주는 대신, 그 계약을 외톨이 점(화살표가 하나도 없는 점)으로 남겨 두고 F1--F6을 되풀이해요(\ref{sec:rank-divergence}절). 두 번째 분석에서 바뀌는 것은 한 층만 걷는 두 걷기뿐이에요. 모든 스펜더는 허락 그래프 안에 있고, 모든 트레이더는 C-PR이 걷는 합집합 그래프 안에 있기 때문이에요. W1의 다른 대비는 라벨을 본 뒤에 계산했고, 결정과 따로 보고해요. 각 페이지랭크와 그 자신의 차수를 견준 대비, 그리고 C-PR과 다른 모든 점수를 견준 대비도 여기에 들어요.

\dissertationSubsubheading[sub:neutral-label-protocol]{두 종류의 화살표 어느 쪽으로도 만들지 않은 결과}

GMX 트레이더 라벨은 두 종류의 화살표 어느 쪽으로도 만들지 않았어요. 이 라벨에서 두 화살표 종류를 견주는 비교를, 규칙~B와 같은 커밋에서 등록했어요. 사전 등록(preregistration)에서처럼, 이 비교의 대비, 규칙, 문구는 이 연구가 7월부터 9월까지의 로그를 뽑기 전에, 그리고 W0 청산 라벨을 하나라도 계산하기 전에 적어 두었어요\parencite{nosek2018preregistration}. 고친 점수에는 이것들을 재분석으로 적용해요. 이 비교는 \ref{sub:trader-holdout}절과 \ref{sub:fresh-holdout}절의 트레이더 코호트와 라벨 코드를 두 기간에서 다시 써요. 하나는 W1(주 기간)이에요. 2026년 6월 30일에 동결했고, 라벨은 7월부터 9월까지이며, 트레이더는 $70$개예요. 다른 하나는 W0(보조 기간)이에요. 2026년 3월 31일에 동결했고, 라벨은 4월부터 6월까지이며, 트레이더는 $87$개예요.

허락과 송금을 맞세우는 쌍은 네 개예요. 첫째는 받은 허락 수와 송금 입차수예요. 둘째는 AWP의 활동 재시작을 쓴 엔도스랭크와 AWP예요. 이것은 같은 조건끼리 견주는 페이지랭크 쌍이에요. 셋째는 엔도스랭크와 AWP예요. 넷째는 금액으로 무게를 단 두 층인 C-PR ($\lambda=1$)과 C-PR ($\lambda=0$)이에요. C-PR ($\lambda=0$) 빼기 AWP는 AWP의 무게와 다시 출발이 만드는 차이만 따로 떼어 보여 줘요. 트레이더 라벨 여섯 개를 모두 보고해요. 청산 없이 닫은 비율(주 라벨), 청산이 없었는지 나타내는 표시, 이익을 낸 닫기, 실현 이익, 그리고 닫기 가운데 이익을 낸 비율과 손실이 아닌 비율이에요. 결정은 W1의 청산 없이 닫은 비율에 대한 대비 두 개로 내려요. P는 받은 허락 수 빼기 송금 입차수이고, Q는 활동 재시작을 쓴 엔도스랭크 빼기 AWP예요. 둘 다 짝지은 다시 뽑기 $B=2{,}000$번(시드 42)으로 구한 본페로니 97.5\% 구간(여러 번 비교할 때 넓힌 신뢰구간)을 써요. 닫기가 많을수록 청산이 한 번이라도 일어날 가능성이 커져요. 그래서 $\tau_b$를 라벨 기간의 닫기 횟수로 나눈 묶음(층) 안에서도 계산해요. 묶음은 3--4번, 5--9번, 10--24번, 25번 이상이에요. 그리고 각 묶음의 계약 쌍 수로 무게를 달아 합쳐요. 허락을 받은 트레이더(동결일에 받은 허락 수가 0보다 큰 트레이더)는 모든 라벨에서 다른 트레이더와 견주어요. 또 영향 점검(몇몇 값이 결과를 좌우하는지 보는 점검)에서는 받은 허락 수가 가장 큰, 허락을 받은 트레이더 10개를 빼요. W1에서 허락을 받은 트레이더는 모두, 그리고 W0에서 허락을 받은 트레이더 39개 가운데 38개는 똑같은 코호트 계약 하나에게서 허락을 받았어요. 그래서 허락을 받은 트레이더들은 한 덩어리를 이뤄요. 두 트레이더 코호트에는 같은 계약 56개가 함께 들어 있고, W1에서 허락을 받은 트레이더는 모두 그 안에 있어요. 그래서 W0은 따로 떨어진(독립적인) 두 번째 표본이 아니에요.

R1은 청산을 피하는 데에서 그냥 센 수의 수준으로 앞서는지를 봐요. R1이 성립하려면 W1에서 P의 97.5\% 구간과 층화한 P의 95\% 구간이 0보다 위에 있어야 하고, W0에서 P가 0보다 커야 해요. R2는 페이지랭크 수준에서 앞서는지를 봐요. R2가 성립하려면 W1에서 Q의 97.5\% 구간이 0보다 위에 있어야 하고, W0에서 Q가 0보다 커야 해요. R3은 W1에서 이익을 낸 비율이나 손실이 아닌 비율에 대한 P가 0보다 작고, 그 95\% 구간에 0이 들어 있지 않을 때 적용돼요. 그러면 허락 신호는 청산을 피하는 데에만 들어맞는 신호로 읽어요. 결과마다 결론을 어떤 문장으로 쓸지도 같은 커밋에서 정했어요.

\dissertationSubheading[sec:proxy-evaluation]{구성개념 점검의 설계}

모든 점수가 같은 기간 평가를 받아요. 엔도스랭크, AWP, $\lambda$의 다섯 값에서 계산한 C-PR, 그리고 S-PR이에요. 처리 과정은 점수 $M$과 점검 $q$마다 $\mathbf{s}_M$과 $\mathbf{p}_q$ 사이의 스피어먼의 로($\rho$, 두 줄 세우기가 얼마나 비슷한지 나타내는 수)와 켄달의 타우($\tau$, 두 줄 세우기가 얼마나 같은 순서인지 나타내는 수)를 내놓아요. 이때 95\% 부트스트랩 구간도 함께 내놓아요(\ref{sec:statistical-tests}절). 점검은 모두 값이 클수록 기대되는 평판이 높다는 뜻이 되도록 방향을 맞추었어요. AWP 평가 절차에서와 같은 방식이에요\parencite{do2023awp}.

같은 기간 평가는 같은 시기에 잰 두 지표 가족을 중심으로 짜여 있어요.
\begin{enumerate}
\item 송금(AWP를 평가하는 데 쓴 대리 지표\parencite{do2023awp}): 계약이 받는 송금 활동이에요.
\item 허락: 계약이 스펜더가 되어 받는 허락 활동이에요.
\end{enumerate}
두 가족은 모두 그래프와 가까운 사이예요. 두 가족이 쓰는 재료가 $G_{\text{transfer}}$와 $G_{\text{approve}}$를 이루는 바로 그 이벤트 종류이기 때문이에요. 다만 각 점검은 그래프 전체를 보는 점수가 아니라, 한 계약 주변만 보는 통계예요. AWP와 송금 점검 사이의 일치도, 또는 엔도스랭크와 허락 점검 사이의 일치도는 의도한 구성개념 점검이에요. 이것은 한 구성개념 안에서 앞뒤가 맞는다는 것을 보여 줘요. 그리고 이것은 다른 것과 상관없이 따로 미리 맞히는 힘(독립적인 예측 능력)의 증거가 아니에요. 세 번째 가족은 \ref{sub:proxy-sybil}절의 시빌 보정 활동 요약이에요. 이 가족은 규칙~A의 보조 부분에 들어가지만, 어떤 가설도 이 가족에 기대지 않아요.

기간 밖 점검은 \ref{sec:holdout-protocol}절의 시간 홀드아웃이 맡아요. \textcite{packin2024credit}에서는 탈중앙 신용 점수가 속이 들여다보이지 않게(불투명하게) 될 수 있다고 경고해요. 이 연구에서는 식을 분명하게 공개하고(\ref{sec:proxy-formulas}절), 같은 기간 점검을 시간으로 나눈 시험과 따로 두어서 그런 불투명함을 줄여요.

\dissertationSubheading[sec:proxy-formulas]{대리 지표를 계산하는 정의}

아래에서 $\mathcal{W}$는 매칭 코호트를 뜻해요. 기호는 공개된 평가 처리 과정에서 점검을 계산하는 모듈(코드 묶음)이 쓰는 기호와 같아요.

\dissertationSubsubheading[sub:proxy-transfer]{송금 가족}

기간 안에 계약 $w$가 받은 송금의 집합을 $\mathcal{T}_{\mathrm{in}}(w)$라고 써요. 송금 $e$를 보낸 쪽은 $u_e$라고 써요. 자기 자신에게 보낸 송금도 $\mathcal{T}_{\mathrm{in}}(w)$에 들어가요. 그래서 자기에게 토큰을 보내는 계약은 자기 자신을 보낸 쪽 하나로 세요. 금액 $x_e$는 다섯 토큰에 걸쳐 달러로 더해요. 이제 다음과 같이 정해요.
\begin{align}
\mathrm{in\_degree}(w)
&=
\bigl|\{\,u_e : e\in\mathcal{T}_{\mathrm{in}}(w)\,\}\bigr|,
\label{eq:in-degree}\\
\mathrm{in\_value}(w)
&=
\sum_{e\in\mathcal{T}_{\mathrm{in}}(w)} x_e.
\label{eq:in-value}
\end{align}
두 값은 모두 AWP 논문에서 정한 기준선이에요\parencite{do2023awp}. 이 기준선은 경제적으로 중요한 주소일수록 더 많은 서로 다른 상대에게서, 더 큰 총금액의 송금을 받는다는 생각에서 출발해요. 페이지랭크와 달리 두 값은 한 주소 주변만 보고, 보낸 쪽의 평판은 따지지 않아요. 그래서 AWP를 이 두 값과 견주어 보면 한 가지를 알 수 있어요. 퍼뜨림, 최근 것에 더 큰 무게 주기, 활발한 보낸 쪽에서 다시 출발하기가 AWP의 순서를 그냥 들어온 양에서 얼마나 멀리 옮기는지예요.

\dissertationSubsubheading[sub:proxy-allowance]{허락 가족}

허락 쪽에서는 $w$가 스펜더예요(주인 $\to w$). 식~\eqref{eq:approve-weight}의 마지막 허락을 $a^{\star}$라고 하고, 다음과 같이 정해요.
\begin{align}
\mathrm{in\_approve\_degree}(w)
&=
\bigl|\{\,u : a^{\star}(u,w,\cdot)>0\,\}\bigr|,
\label{eq:in-approve-degree}\\
\mathrm{in\_approve\_value}(w)
&=
\sum_{u,\mathrm{token}} \min\bigl(a^{\star}(u,w,\mathrm{token}),U\bigr).
\label{eq:in-approve-value}
\end{align}
송금 기준선처럼, 이 두 값도 여러 단계로 이어지는 보증 사슬을 따라가지 않아요. 들어온 허락을 세거나 더하기만 해요. 받은 허락 금액에는 C-PR의 상한 $U$를 써요.

\dissertationSubsubheading[sub:proxy-sybil]{시빌 보정 안정성 가족}

들어온 것을 세는 값들은 자기 자신에게 보낸 송금($u=w$)을 빼요. $N_{\mathrm{tx}}(w)$는 $w$가 다른 주소에게서 받은 송금의 수예요. $N_{\mathrm{uniq}}(w)$는 송금을 보내온 서로 다른 상대의 수예요. 활동 기간과 활동한 달 수는 또 다른 행(기록 한 줄) 묶음을 써요. 코호트 계약이 끼어 있는 송금은, 보낸 쪽이 코호트에 있으면 보낸 쪽의 것으로 쳐요. 그렇지 않으면 받는 쪽의 것으로 쳐요. 그래서 코호트 계약 두 개 사이의 송금은 보낸 쪽에만 세어요. $t_{\min}(w)$와 $t_{\max}(w)$는 $w$의 것으로 친 행들 가운데 가장 이른 시각과 가장 늦은 시각이에요. $N_{\mathrm{mon}}(w)$는 서로 다른 달력 달을 센 두 수 가운데 큰 쪽이에요. 하나는 $w$가 다른 주소에게서 송금을 받은 달의 수이고, 다른 하나는 $w$의 것으로 친 행들이 있는 달의 수예요. 이제 다음과 같이 정해요.
\begin{align}
\mathrm{inbound\_counterparty\_ratio}(w)
&=
\frac{N_{\mathrm{uniq}}(w)}{\max(N_{\mathrm{tx}}(w),1)},
\label{eq:counterparty-ratio}\\
\mathrm{transfer\_tenure\_days}(w)
&=
\frac{t_{\max}(w)-t_{\min}(w)}{86400}\quad\text{(초를 날로 바꿈)},
\label{eq:tenure}\\
\mathrm{active\_months}(w)
&=
N_{\mathrm{mon}}(w).
\label{eq:active-months}
\end{align}
상대 비율은 적은 주소에게서 송금을 많이 받는 계약에 낮은 값을 줘요. 이것은 간단한 자전 거래(자기끼리 사고팔아 거래가 많은 척하기) 모양이에요. 하지만 흔한 시빌 모양에는 낮은 값을 주지 못해요. 흔한 시빌 모양은 새로 만든 주소 여러 개가 송금을 하나씩 보내는 거예요. 이렇게 송금을 받은 계약은 가장 큰 값인 1에 가까운 값을 받아요. 활동 기간과 활동한 달 수는 점수 기간 동안 온체인 활동이 얼마나 꾸준히 이어졌는지를 봐요. 그래서 이 대리 지표들은 약한 시빌 점검이에요. 짜고 움직이는 공격자들은 세 가지를 모두 채울 수 있어요.

\dissertationSubheading[sec:statistical-tests]{통계 검정}

평판 점수도 대리 지표도 빚을 갚았는지 못 갚았는지처럼 두 가지 값만 갖는 표시가 아니에요. 둘 다 이어지는 값이고, 순위로 다뤄요. AWP 평가 절차와 \textcite{cronbach1955construct}의 구성 타당도 접근을 따라서, 순서가 얼마나 맞는지는 순위 상관으로 재요\parencite{do2023awp}. 순위 상관은 어느 쪽 변수든 순서를 바꾸지 않는 방식으로 크기를 바꾸면 달라지지 않아요.

\begin{itemize}
\item 스피어먼의 로 $\rho$: 계약 $i$의 원래 점수와 원래 대리 지표 값에 매긴 중간 순위 $R_i$와 $S_i$ 사이의 피어슨 상관계수(두 값이 함께 커지고 작아지는 정도를 재는 수)예요. 동점인 값들은 자기들이 차지한 자리 번호의 평균을 나누어 가져요(\ref{sec:rank-correlation-theory}절). 이 값은 한 방향으로 꾸준히 함께 움직이는 정도를 재요. 평판이 높을수록 대리 지표 값도 꾸준히 높으면, 원래 점수가 곧은 선에서 얼마나 벗어나 있든 $\rho$는 양수가 돼요.

\item 켄달의 타우 $\tau$: 보고한 값은 모두 식~\eqref{eq:kendall}의 동점 보정 $\tau_b$이고, 원래 점수와 원래 대리 지표로 계산했어요. 동점이 없으면 $(1+\tau)/2$는 아무렇게나 고른 쌍 하나를 두 순위가 같은 방향으로 줄 세울 확률이에요. 동점이 있으면 이렇게 읽는 방법은 $\tau_b$가 아니라 $\tau_a=(n_c-n_d)/n_0$에 맞아요. 매칭 코호트의 페이지랭크 점수에서는 동점이 드물어요(\ref{sec:rank-divergence}절). 하지만 원시 차수와 라벨에서는 동점이 흔해요. 짝지은 대비의 두 쪽은 동점으로 묶는 쌍의 비율이 크게 다를 수 있어요. 그냥 센 수와 페이지랭크가 그래요. 이때 $\tau_b$는 동점이 더 많은 쪽에 더 작은 분모를 줘요. 그래서 그 대비는 $\tau_b$에서는 동점이 많은 쪽에 유리하게 나오고, $\tau_a$에서는 그렇지 않을 수 있어요. 차수 대비와 P는 $\tau_b$로만 계산했기 때문에, 이 문제를 그대로 안고 있어요.

\item 둘 다 쓰는 까닭: AWP 평가 방식도 둘 다 보고해요\parencite{do2023awp}. 둘 다 보고하면 결론이 계수 하나에만 매달리지 않아요. 점수 분포는 꼬리가 두꺼워요(아주 큰 값이 드물게 나와요). 대부분의 계약은 0 가까이에 있고, 몇 안 되는 허브(화살표가 많이 모이는 점)가 전체 몫의 많은 부분을 차지해요. 그래서 원래 점수로 계산한 피어슨 상관계수 대신 순위 계수를 써요. 피어슨 상관계수는 허브들이 좌우해 버리기 때문이에요.

\item 가족 묶기: 한 가족은 그 안 대리 지표들의 단순 평균으로 요약해요. 예를 들어 송금 가족 평균 $\tau$는 $\tfrac{1}{2}(\tau_{\mathrm{in\_degree}}+\tau_{\mathrm{in\_value}})$예요. 결론은 신뢰구간(믿을 만한 범위)에서 끌어내요. 정해 둔 $\tau$ 값을 넘는지를 판단 기준으로 쓰지 않아요.

\item 불확실성: 표본을 뽑는 데서 오는 불확실성은 계약 단위의 짝지은 부트스트랩으로 어림해요\parencite{efron1979bootstrap}. 평가하는 코호트를, 뽑은 것을 되돌려 넣으며 $B=400$번 다시 뽑아요(시드 42, 설정 파일에 고정). 그리고 번호 행렬 하나를 모든 방법과 모든 대리 지표에 다시 써요. 그래서 다시 뽑아 구한 통계는 방법끼리 짝지어져요. 대리 지표마다의 $\tau$, 지표 가족 평균, 방법 사이의 차이마다 부트스트랩 분포의 2.5번째 백분위와 97.5번째 백분위가 95\% 백분위 구간을 이뤄요\parencite{diciccio1996bootstrap}. $\tau_b$는 U-통계량(쌍을 모두 견주어 평균 낸 값)의 매끄러운 함수예요. 그래서 계약들이 서로 독립이라면, 매칭 코호트와 스펜더 코호트의 크기에서 백분위 구간은 1차 정확도를 가져요. 하지만 주인을 함께 가지거나 한 서비스(프로토콜)에 속한 계약들은 서로 독립이 아니에요. 그래서 구간이 너무 좁을 수 있어요. 다시 뽑는 횟수가 적어서, 동점 보정 $\tau$의 $O(B\,n\log n)$ 비용을 모든 방법과 대리 지표에 대해 감당할 수 있어요. 그 대가로 구간 끝마다 몬테카를로 오차(무작위로 다시 뽑아서 생기는 흔들림)가 생기고, 그 크기는 부트스트랩 표준 오차의 약 $0.13$배예요. \ref{sub:neutral-label-protocol}절의 분석만 이와 달리 $B=2{,}000$을 써요.

\item 차이 검정: 모든 비교는 같은 다시 뽑은 표본 위에서 하는 짝지은 대비예요. 다시 뽑은 표본 하나하나마다 $\Delta\tau=\tau_A-\tau_B$를 계산하고, 그 부트스트랩 분포의 95\% 백분위 구간을 보고해요. 구간이 0을 벗어나면 그 대비는 차이가 있다고 쳐요. 라벨을 본 뒤에 계산한 대비는 백 개가 넘는데, 그 구간들은 여러 번 비교한 것에 맞추어 넓히지 않았어요. 그러니 그 가운데 몇 개는 우연만으로도 0을 벗어날 거예요. 그래서 한 기간에서만 차이가 드러난 사후 대비는 가설로 읽어요. 두 기간 모두에서 같은 방향으로 차이가 드러난 대비는 우연에 덜 휘둘린 것으로 읽어요. 다시 뽑기가 짝지어져 있어서, 대비는 두 추정값의 공분산(두 추정값이 함께 흔들리는 정도)을 그대로 지녀요. 곧 $\operatorname{Var}(\hat\tau_A-\hat\tau_B)=\operatorname{Var}\hat\tau_A+\operatorname{Var}\hat\tau_B-2\operatorname{Cov}(\hat\tau_A,\hat\tau_B)$이에요. 그래서 공분산이 양수이면, 대비의 구간은 두 구간을 따로 볼 때 짐작되는 것보다 좁아요. 규칙을 판정하는 대비는 이 연구가 판정에 쓸 자료를 뽑기 전에 정해 두었어요. 규칙~A의 대비(\ref{sec:holdout-protocol}절)는 로그를 하나도 뽑기 전에 정했어요. 규칙~B의 대비와 \ref{sub:neutral-label-protocol}절의 대비는 W1 로그를 뽑기 전에 정했어요. 수정 분석은 등록 분석의 결과를 안 뒤에 이 대비들을 써요. W0의 새 허락 쌍에서 C-PR을 송금 층만 걷는 것과 견준 대비는 여기에 들지 않았어요. 이 대비는 나올 때마다 탐색적이라고 표시해요. 같은 기간에서는 여덟 개 대비가 C-PR과 S-PR을 각 가족에서 앞선 층(한 층만 걷는 것)과 견줘요. 송금 가족에서는 송금 걷기, 허락 가족에서는 허락 걷기, 시빌 안정성 가족에서는 두 층 모두예요. 그리고 여섯 개 대비가 C-PR을 송금 가족과 허락 가족에서는 AWP와, 허락 가족과 송금 가족에서는 엔도스랭크와, 시빌 안정성 가족에서는 둘 다와 견줘요. 이 열네 개 가운데 어느 것도 계획에서 정해 두지 않았어요. 그리고 그 가운데 두 개, 곧 허락 가족에서 C-PR과 AWP를 견준 것과 송금 가족에서 C-PR과 엔도스랭크를 견준 것은 저장해 둔 점수로 나중에 계산했어요. 기간 밖에서는 C-PR과 다른 모든 점수, 그리고 C-PR과 두 원시 차수를 견준 대비(표~\ref{tab:cpr-contrasts})도 라벨을 본 뒤에 계산했어요. 엔도스랭크와 AWP가 들어가는 같은 기간 대비 다섯 개도 미리 정해 두지 않았고, 부록~\ref{ch:appendix-er-awp}에 보고해요.

\item 방향과 빈칸 채우기: 대리 지표는 모두 값이 클수록 평판이 좋다는 뜻이 되도록 방향을 맞춰요. 평가하는 코호트에 맞추어 다시 줄을 맞춘 뒤, 어떤 대리 지표 자료에 없는 계약은 0으로 채워요. 거래가 듬성듬성하므로 이것은 조심스러운(보수적인) 선택이에요. 그리고 두 점수 방식 모두 똑같이 이렇게 해요. 이 단계를 거치면 평가하는 계약마다 모든 점수 벡터와 대리 지표 벡터에 값이 있어요. 그래서 상관은 모두 코호트 전체를 써서 계산해요.

\item 비열등성(정한 폭보다 더 나쁘지 않음): 규칙~B의 세 조건 가운데 두 개는 C-PR이 송금 층만 걷는 것보다 미리 정한 허용 폭 $\delta=0.01$보다 더 많이 뒤처지지만 않으면 돼요. 이런 대비는 95\% 구간의 아래 끝이 $-\delta$보다 위에 있으면 통과해요. 이 검정은 차이 검정보다 약하고, 그 라벨에서 더 낫다는 것(우월성)에 대해서는 아무것도 말해 주지 않아요. 이 검정은 C-PR이 중요한 것을 하나도 잃지 않는다는 주장을 시험할 때 써요.

\item 하지 않는 것: 이 연구는 인과 추론(무엇이 원인이고 무엇이 결과인지 밝히기)을 하지 않아요. 훈련 자료와 시험 자료를 나누어 지도 학습 모델(정답을 보고 배우는 모델)을 훈련하지도 않아요. $p$값도 보고하지 않아요. 대신 대비의 구간이 영가설 값(차이가 없다는 값)을 벗어나는지를 물어요.

\end{itemize}

\dissertationSubheading[sec:computational-benchmark]{계산 벤치마크(주장~C5)}

주장~C5는 점수마다 계산하는 데 드는 비용에 대한 거예요. 벤치마크는 화살표 목록을 점수로 바꾸는 풀이 프로그램 호출의 시간을 재요. 화살표 목록을 만드는 데 한 번 드는 비용은 빼요. 이 비용은 대부분 불러오기와 미리 손질하기예요. 점수를 주기적으로 새로 고치는 점수 서비스라면 이 비용을 여러 번에 나누어 부담할 거예요.

\begin{itemize}
\item 시간을 재는 부분: 화살표 목록과 AWP의 다시 출발 무게는 시계를 켜기 전에 만들어요. 시간을 재는 호출은 점에 번호를 매기고, 무게를 단 듬성듬성한 인접 행렬(어느 점에서 어느 점으로 화살표가 있는지 적은 표)을 짜 맞추고, 줄마다 합이 1이 되게 맞추고, 재시작 벡터를 만들고, $d=0.85$, $\varepsilon=10^{-8}$, $I_{\max}=300$으로 거듭제곱 반복(같은 계산을 되풀이하기)을 돌려요. 행렬 짜 맞추기와 반복 한 바퀴는 둘 다 $|E|$와 $|V|$에 따라 늘어나는 시간이 걸려요. 실제로 운영한다면 업데이트할 때마다 둘 다 다시 해야 해요. 일치도 점수도 같은 그래프에서 나와요.

\item 재는 방법: 모든 설정은 한 세션 안에서 재요. 먼저 시간을 재지 않는 준비 실행을 한 번 해요. 가장 많이 쓴 메모리는 이 준비 실행에서만 추적해요. 그다음 시간 잰 실행을 다섯 번 해요. 이 다섯 번에서 실행 시간의 평균과 표준편차(SD, 값들이 얼마나 흩어져 있는지), 반복 횟수를 보고하고, 입력 그래프의 점 수와 화살표 수도 함께 보고해요. 같은 설정의 같은 그래프가 세 곳에 나와요. 주 벤치마크, 감쇠 계수를 바꿔 보는 실험의 $d=0.85$ 행, 그리고 $n=6{,}483$인 규모 단계예요. 이 마지막 단계가 코호트 그래프 전체예요(\ref{sub:expanded-pool}절). 이 표들은 일부러 측정 하나를 다시 써요. 그래서 이것들은 한 양을 따로따로 잰 측정이 아니에요. 실행 시간 표에는 모두 화살표 수를 적어요. 반복 한 바퀴에 $O(|E|)$만큼 일하는 알고리즘에서는 화살표 수가 비용을 정하기 때문이에요\parencite{langville2006pagerank}. 실행 시간은 기계가 읽을 수 있는 출력에 기록해요(\ref{sec:reproducibility}절).

\item 코호트: 주 벤치마크는 매칭 계약 코호트($n=6{,}483$)의 그래프에서 돌려요. 그래프 크기에 따라 실행 시간이 어떻게 늘어나는지 보려고, 코호트에서 정해진 방식으로 부분 표본을 뽑아요(\ref{sec:scaling-robustness}절).

\item 공정함: 모든 알고리즘은 같은 풀이 프로그램, 허용 오차, 반복 상한, 감쇠 계수(감쇠 계수를 바꿔 보는 실험은 빼고), 씨앗 집합으로 돌려요. 그래서 비용 차이는 구현에서 오지 않고, 그래프와 AWP의 재시작 벡터에서 와요. 이 재시작 벡터는 반복 횟수를 바꿀 수 있어요. 하드웨어와 소프트웨어는 벤치마크 요약에 적혀 있어요. 두 점수의 비용은 시간 비율로 견주고, 시간 비율마다 그에 맞는 화살표 수 비율을 옆에 함께 보고해요.

\end{itemize}

\dissertationSubheading[sec:scaling-robustness]{규모와 강건성을 보는 추가 분석}

이 연구에서는 그래프 크기에 따라 실행 시간이 어떻게 늘어나는지도 살펴봐요. 또 감쇠 계수, 금액 무게의 기울기, 표본을 바꾸어도 가족 수준의 모양이 그대로 남는지도 살펴봐요. 이 분석들은 주장~C5에 조건을 달아 주고, 가족 수준 결론이 어디까지 통하는지 시험해요. 어느 분석도 주 표본을 대신할 다른 표본을 내놓지 않아요. 따로 말하지 않으면, \ref{ch:results}장의 일치도 주장은 매칭 코호트에 대한 거예요.

\dissertationSubsubheading[sub:expanded-pool]{실행 시간의 규모 변화를 보기 위한 코호트 부분 표본}

규모 단계는 시드 문자열 \texttt{contract-benchmark-v1}을 쓰는 SHA256 부분 표본 뽑기로 매칭 코호트에서 뽑아요. 이 뽑기는 늘 같은 결과를 내요. 계약들은 아래 값의 순서로 줄을 세워요.
\begin{center}
\texttt{SHA256(}\textit{시드}\texttt{:}\textit{주소}\texttt{)}, 여기서 \textit{시드} $=$ \texttt{contract-benchmark-v1},
\end{center}
곧, 시드와 쌍점(:)과 소문자로 바꾼 주소를 이어 붙인 문자열 하나의 해시값으로 줄을 세워요. 부분 표본은 이 순서에서 맨 앞의 계약 $n$개예요. 시드 문자열은 이 연구의 설정 파일에 고정한 상수예요. 보고하는 단계는 네 개예요. $n=1{,}000$, $2{,}000$, $5{,}000$, 그리고 코호트 전체인 $n=6{,}483$이에요. 해시 계산은 언제나 같은 결과를 내요. 그래서 이 부분 표본은 의사 난수(컴퓨터가 계산으로 만드는 가짜 무작위 수)로 뽑을 때처럼 라이브러리 버전에 따라 달라지지 않아요. 그리고 각 단계에는 바로 앞 단계가 모두 들어 있어요.

이 설계의 두 가지 특징 때문에, 이 실험이 보여 줄 수 있는 것에는 한계가 있어요. 첫째, 단계마다 그래프는 $T_{\mathrm{obs}}$ 시점 코호트 그래프의 화살표 가운데, 양 끝 중 적어도 하나가 뽑힌 계약인 화살표를 남겨요. 풀이 프로그램의 점 집합은 그 화살표들의 양 끝 점으로 이루어져요. 그래서 화살표 수 $|E|$는 단계가 커질수록 늘어나지만, $n$에 비례해서 늘지는 않아요. 상대가 많은 계약은 화살표를 많이 데려오기 때문이에요. 다른 하나는 어느 단계도 코호트 그래프 전체보다 크지 않다는 점이에요. 그래서 이 실험은 그 크기까지 실행 시간이 $|E|$에 따라 어떻게 늘어나는지만 보여 주고, 더 큰 그래프에 대해서는 아무것도 말해 주지 않아요. 규모 표는 모두 두 방법의 단계별 $|V|$와 $|E|$를 보고해요. \ref{ch:results}장에서 규모에 대해 하는 말은 그 열들이 보여 주는 범위를 넘지 않아요.

\dissertationSubsubheading[sub:damping-protocol]{감쇠 계수 민감도}

매칭 코호트에서 그래프와 대리 지표는 그대로 두고, 감쇠 계수를 $d\in\{0.75,0.85,0.95\}$로 바꿔 봐요. $d$가 작으면 순간 이동이 늘어나요. 반대로 크면 여러 단계를 거치는 길에 더 큰 무게가 실리고, 수렴(값이 더는 거의 바뀌지 않게 됨)이 느려질 수 있어요. 이 실험은 허락 가족과 송금 가족에서의 일치도가 순간 이동 확률에 따라 달라지는지 물어요.

\dissertationSubsubheading[sub:top-tokens]{금액 기울기 민감도}

이 점검에서는 $V$의 기울기 $b$에 10을 곱하거나 10으로 나누어, $T_{\mathrm{obs}}$에서 엔도스랭크와 AWP를 다시 계산해요. 그러면 무게가 절반이 되는 금액이 중앙값의 10분의 1과 중앙값의 10배로 옮겨 가요. 이 점검은 계획에 있던 20개 토큰 점검을 대신해요. 토큰을 다섯 개만 쓰면 그 점검은 뜻이 없어지기 때문이에요(\ref{sec:amount-revision}절). 이 점검은 금액 무게가 더 가파르게 또는 더 완만하게 올라갈 때 가족 수준 결과가 바뀌는지 물어요.

\dissertationSubsubheading[sub:sample-size-protocol]{표본 정의 민감도}

\ref{sub:expanded-pool}절의 규모 단계마다, 엔도스랭크와 AWP의 지표 가족 평균을 그 단계의 계약에서 다시 계산해요. 점수는 그 단계의 그래프로 매겨요. 단계들은 크기와 그래프 두 가지에서 코호트와 달라요. 단계에서 빠진 계약들의 화살표가 없으니, 계약의 점수가 바뀌어요. 그래서 단계별 계수는 표본 정의에 대한 민감도 점검이지, 안정성의 증거가 아니에요. 이 계수들에는 구간도 계산하지 않았어요. 묻는 것은 표본이 바뀔 때 각 계수가 얼마나 움직이는지예요. 이 분석은 계획에 없었어요.

\dissertationSubheading[sec:reproducibility]{재현성}

\ref{ch:results}장의 표는 하나도 손으로 옮겨 적지 않았어요. 1단계의 표가 디스크에 저장되면, 평가 구동 프로그램이 모든 점검, 구간이 붙은 모든 일치도 통계, 모든 홀드아웃 결과, 벤치마크를 다시 계산해요. 그리고 그것들을 기계가 읽을 수 있는 요약 파일에 적어요. 따로 떨어진 단계 하나가 이 요약 파일들로 \ref{ch:results}장과 부록~\ref{ch:appendix-eval}의 표를 내보내요. 버전을 매긴 설정 파일 하나에 다음이 모두 들어 있어요. 기간, 코호트 기준값, 이벤트 토픽 ID, 감쇠 계수, 로지스틱 감쇠의 매개변수, $\lambda$ 값, 부트스트랩 횟수와 시드, 벤치마크 시드, 비용 예산이에요. 수정 분석은 여기에 토큰 목록, 하루 가격, \ref{sec:amount-revision}절의 상수들, 그리고 각 점수 기간의 시작을 더해요.

계획, W1의 등록 문서, 수정 분석의 문서는 공개 저장소 \url{https://github.com/abitocodes/phd_works}에 있어요. 계획과 등록 문서는 커밋 \texttt{43d8045}와 \texttt{0d3a057}에 있어요. 수정 분석은 \texttt{docs/revision\_price\_weighting.md}와 \texttt{docs/revision\_one\_year\_window.md}에 있어요. 다시 똑같이 해 볼 수 있게 해 주는 요소는 다음과 같아요. SHA256 해시로 코호트에서 정해진 방식으로 부분 표본 뽑기, 고정한 페이지랭크 허용 오차와 반복 상한, 코호트마다 부트스트랩 시드 하나와 다시 뽑기 번호 행렬 하나, \ref{sec:proxy-formulas}절의 대리 지표 식, 모든 실행의 시간을 출력에 남기는 \ref{sec:computational-benchmark}절의 벤치마크 절차, 결합 무게 $\lambda$와 \ref{sec:holdout-protocol}절의 홀드아웃 논리(둘 다 이 연구가 로그를 하나도 뽑기 전에 설정 파일에 고정), 정해 둔 블록에서 읽은 계정 종류, 그리고 스캔~B의 추출 기록 파일이에요. 모든 표의 바탕이 되는 기계가 읽을 수 있는 요약 파일은 설정 파일, 계획, 등록 문서와 함께 Git 저장소의 \texttt{experiments/data/3-published-results-for-thesis/} 아래에 커밋되어 있어요(수정 분석은 \texttt{revised-usd-12m/}). 이 폴더는 소스 코드와 같은 폴더 안에 있어요. 부록~\ref{ch:appendix-eval}에 GitHub 저장소 URL, 이 논문의 표를 만든 정확한 커밋 해시, 설정 파일, 요약 파일이 정리되어 있어요.

\dissertationSubheading[sec:sampling-bias]{표본에 대해 생각할 점과 빠진 자료}

매칭 코호트는 아비트럼의 모든 주소에서 무작위로 뽑은 표본이 아니라, 규칙을 정해 만든 표본이에요. 코호트에 들려면 동결일 전 12개월 동안 서로 다른 주인 적어도 세 명에게서 다섯 토큰 가운데 하나에 대해 0이 아닌 허락을 받아야 해요. 송금 활동은 여기에 아무 역할도 하지 않아요. 그래서 같은 기간 결과는 이 몇 달 동안 다섯 토큰에 대해 적어도 세 주인에게서 허락을 받은 아비트럼 원의 계약에 대한 거예요. 홀드아웃 결과는 코호트의 허락 자료에 나오는 스펜더 계약에 대한 거예요. EOA와, 이 토큰들에 대해 허락한 주인이 세 명보다 적은 계약은 처음부터 빠지도록 설계했어요(\ref{sec:data-sources}절).

자료는 몇 가지 방식으로 빠질 수 있어요. 읽어 낼 수 없는 금액은 0으로 읽어요. 그러면 송금은 버려지고, 허락은 0으로 되돌아가요. 엔도스랭크는 표준 ERC-20 이벤트를 남기지 않는 허락을 보지 못해요. AWP는 오프체인에서 한 송금이나, 수탁 계정(거래소 같은 곳이 대신 맡아 관리하는 계정) 안에서 한 송금을 보지 못해요. 그래서 두 점수 모두 온체인에서 표준 토큰으로 하는 행동 쪽으로 치우쳐 있어요. 이 치우침은 온체인 평판 연구로서는 받아들일 만해요. 하지만 결과를 한 서비스의 모든 계약으로 넓혀 적용할 때는 중요해져요.

\dissertationSubheading[sec:implementation-notes]{구현에 대한 메모}

페이지랭크 코드는 파이썬으로 썼어요. SciPy의 듬성듬성한 이동 연산자와 NumPy의 거듭제곱 반복을 써요. 벤치마크를 돌린 기계는 부록~\ref{sec:pipeline-stages}에 적혀 있어요. 무작위성은 부트스트랩에서만 들어오고, 그 시드는 고정되어 있어요. 내려받은 표와 설정 파일이 있으면 평가 결과는 언제나 똑같이 나와요. 플랫폼마다 소수 계산에서 아주 작은 차이가 생길 것으로 예상해요. 하지만 그 차이는 소수점 아래 세 자리까지 나타내는 $\tau$ 지표의 정밀도에 비하면 무시할 만해요.

\dissertationSubheading[sec:ethics]{윤리}

이 연구는 누구나 볼 수 있는 블록체인 자료만 써요. 그리고 사람을 연구 대상으로 삼지 않아요. 이 연구는 기관의 윤리 위원회(Unisa SoC/CAES)의 검토를 받았어요. 부록~\ref{ch:appendix-ethics}에 윤리 승인(연구가 윤리 기준을 지킨다고 인정받음) 증명서가 들어 있어요(참조 번호 [Unisa SoC/CAES 윤리 승인 번호: TBD], 날짜 [TBD]).

\begin{itemize}
\item\phantomsection\label{sub:privacy-anonymity}
  사생활과 익명성: 아비트럼 원의 주소는 가명(진짜 이름 대신 쓰는 이름)이에요. 주소만으로는 그 뒤에 있는 실제 세상의 신원이 드러나지 않아요\parencite{packin2024credit}. 하지만 가명성(진짜 이름 대신 주소로만 활동함)은 익명성(누구인지 전혀 알 수 없음)과 같지 않아요. 이론적으로는, 마음먹은 공격자가 거래소 입금, 온라인 프로필, 어림 규칙을 써서 주소를 한 사람과 이을 수도 있어요. 이 연구에서는 일부러 그런 연결을 하나도 하지 않아요. 점수를 매긴 주소는 계약이고, 어떤 EOA에 대해서도 결과를 보고하지 않아요. 이 연구의 자료 수집은 온체인 \texttt{Approval}, \texttt{Transfer}, GMX~V2 포지션 이벤트로 한정돼요. 오프체인 개인 식별 정보(IP 주소, 거래소 KYC(신원 확인) 기록, 소셜 미디어 계정 이름, 익명성을 벗기는 그래프 같은 것)는 긁어모으지 않았어요. 주소들을 실제 세상의 주체로 묶지 않았어요. 개인을 알아볼 수 있을 것 같은 주소의 순위표를 공개하지도 않았어요. 신원을 알아내려고 한 주소의 거래 기록을 다른 데이터베이스와 합쳐 보지도 않았어요. 공개한 결과는 모아서 낸 값(상관계수, 실행 시간, 지표 가족 평균)뿐이고, 이름이 밝혀진 참여자의 개별 순위는 들어 있지 않아요.

\item\phantomsection\label{sub:data-minimization}
  자료 최소화: 스캔~A는 2023년 10월 1일부터 2026년 6월 30일까지의 모든 ERC-20 \texttt{Approval}과 \texttt{Transfer} 이벤트를 이 연구의 클라우드 프로젝트 안에 보관해요. 코호트 규칙에 필요하기 때문이에요. 그 뒤의 모든 질의와 내려받은 모든 표는 코호트의 에고 네트워크로, 그리고 그래프 만들기와 점검에 필요한 칸으로만 한정돼요. 그 칸은 관련된 주소, 토큰 금액, 시각이에요. 뽑은 표는 연구 처리 과정을 감사(꼼꼼히 확인)할 수 있도록 보관해요. 사용자의 익명성을 벗기려는 목적으로는 공개하지 않을 거예요.


\item\phantomsection\label{sub:data-usage-compliance}
  자료 사용과 규정 지키기: 모든 질의는 \texttt{bigquery-public-data.goog\_blockchain\_arbitrum\_one\_us} 데이터셋에 올라 있는 공개 자료를 대상으로 돌렸어요. 그리고 구글 클라우드의 이용 약관과 요청 횟수 제한을 지켰어요. 미리 돌려 보기와 연구 예산이 쓸데없는 스캔을 막았어요. 계정 종류와 토큰 자료는 읽기만 하는 RPC 호출로 읽었어요. 하루 가격은 공개 DefiLlama API에서 가져왔어요. 이 연구는 거래를 하나도 보내지 않았어요. 특별한 권한이 있거나 공개되지 않은 색인 서비스의 접속 주소에도 접속하지 않았어요. 그리고 네트워크의 정상 작동을 방해할 수 있는 방식으로 스마트 계약을 부른 일도 전혀 없어요.


\item\phantomsection\label{sub:fairness-transparency}
  공정함과 투명성: 탈중앙 신용 점수는 불투명함뿐 아니라 공정함에 대해서도 걱정을 낳아요\parencite{packin2024credit}. 그래프에 바탕을 둔 평판 체계는 새 계약, 거의 쓰이지 않는 계약, 공개된 허락을 받지 못한 계약을 모두 낮게 평가할 수 있어요. 엔도스랭크와 AWP 모두 이런 구조적인 치우침을 지니고 있어요. 예를 들어 허락 그래프에서는, 쓸 권한을 한 번도 받지 못한 주소가 받은 주소보다 불리해요. 한편 송금 그래프는 돈이 흐르는 네트워크 안에 깊이 들어 있는 주소를 더 좋게 봐요. \ref{ch:discussion}장에서 한계를 다뤄요. 그 가운데에는 콜드 스타트 효과(새로 생겨 기록이 적은 계약이 불리한 일)와 허락 그래프가 시빌 공격에 약한 점이 있어요.


\item\phantomsection\label{sub:responsible-disclosure}
  책임 있는 공개와 영향: 엔도스랭크는 디파이(DeFi, 탈중앙 금융)에서 위험을 따져 보는 데 도움을 줄 수 있어요. 하지만 금융 조언, 신용 조언, 소비자 신용 보고 상품은 아니에요. 빌려주기, 거래, 규정 준수에 대한 결정에는 여전히 사람의 감독이 필요해요. \ref{ch:discussion}장에서는 공격자가 점수를 조작할 가능성을 인정하고 다뤄요. 예를 들면 $G_{\text{approve}}$에 대한 시빌 공격이나 $G_{\text{transfer}}$에서의 자전 송금이에요. 앞으로 실제로 운영하려면 운영 규칙(거버넌스), 이의 제기 절차, 공정성 감시가 더 필요할 거예요. 이것들은 모두 이 연구의 범위 밖이에요.


\item\phantomsection\label{sub:methodological-commitments}
  방법에서 미리 약속한 것 정리: 결과를 해석하기 전에, 설계는 이미 체인, 기간, 코호트 규칙과 계정 종류 규칙, 그래프 식, 대리 지표 정의, 부트스트랩 방법과 시간 재기 방법, 통계 검정을 정해 두었어요. 그 뒤에 바꾼 것은 하나뿐이고, \ref{sec:amount-revision}절에서 설명해요. 주된 주장은 매칭 계약 코호트와 스펜더 홀드아웃에서 모든 점수를 함께 다뤄요. 그리고 각 페이지랭크를 그것이 매끄럽게 다듬는 원시 차수와 견줘요.
\end{itemize}

\ref{ch:results}장은 이 절차가 아비트럼 원에서 어떤 결과를 내는지 보고해요.
